\documentclass[12pt,openany]{book}

\usepackage[margin=0.9in]{geometry}
\usepackage{amsmath,amssymb,amsthm}
\usepackage{graphicx}
\usepackage[T1]{fontenc}

\setlength{\parindent}{0pt}
\setlength{\parskip}{4pt}

\theoremstyle{definition}
\newtheorem{definition}{Definition}[chapter]

\usepackage[hidelinks]{hyperref}
\hypersetup{
  pdftitle={Linear, Geometric, and Structural Transformations: An Axiomatic and Coordinate-Free Perspective},
  pdfauthor={Jos\'e Portillo},
  pdfsubject={Linear algebra, geometry, tensors, and logarithmic structure},
  pdfkeywords={linear transformations, inner products, tensors, tensor product, logarithms}
}

\begin{document}

\pagestyle{plain}

\frontmatter

\begin{titlepage}
\thispagestyle{empty}

\vspace*{2.2cm}

{\Large\bfseries
Linear, Geometric, and Structural Transformations\par}
\vspace{0.4cm}

\vspace{1.4cm}
{\Large Jos\'e Portillo\par}

\vspace{1.8cm}
\rule{\textwidth}{0.6pt}
\vspace{0.35cm}

{\small
\textbf{Release date:} February 15, 2026\par
\textbf{Version:} 0.4\par
}

\vfill

\end{titlepage}


\thispagestyle{empty}
\null\vfill
\begin{center}
To the memory of\par
Mercelena Portillo Holskin,\par
Eldon C. Hall,\par
Hugh Blair-Smith\par
\vspace{1cm}
To my family
\end{center}
\vfill\null
\clearpage

\tableofcontents

\chapter*{Preface}
\addcontentsline{toc}{chapter}{Preface}

Every algorithm, however sophisticated, reduces to a finite sequence of
elementary operations: additions, multiplications, comparisons.
Every physical system that performs computation --- whether built from vacuum
tubes, transistors, optical waveguides, or any substrate yet to be conceived
--- implements those same operations through different physical means.
What persists (underlies) all such changes of medium and representation is
a layer of mathematical structure.
This monograph is concerned with that layer, which remains invariant.

\medskip

The question that organizes the exposition is therefore \emph{what} is being computed --- and what
structural properties that object possesses independently of the coordinate
system, basis, or hardware in which it happens to be expressed.
This monograph began as an attempt to understand what remains invariant when
formulas, coordinates, and representations are removed.

\medskip

A recurring distinction will be made between an axiomatic theory and a
concrete model of that theory.
In a refined axiomatic treatment of Euclidean geometry, objects such as
\textit{point}, \textit{line}, and \textit{plane} are taken as primitive
(undefined) terms, and the axioms specify only the relations they must satisfy.
A Cartesian coordinate system does not redefine these primitives; rather, it
provides a \emph{model} in which points are represented by tuples
in~$\mathbb{R}^n$, lines by solution sets of linear equations, and incidence
becomes set membership~\cite{iribarren2}.
As stated in~\cite{wilder}, this is \textit{the assignment of meanings}.
This perspective helps place coordinate formulas as representations of
structure, and not as the structure itself.
The same distinction runs through every chapter of this text.

\medskip

Vector spaces are introduced as abstract sets equipped with addition and scalar
multiplication.
Linear transformations are defined as the maps that preserve this structure.

\medskip

By introducing bilinear forms and, ultimately, inner products, vectors are
allowed to interact in a way that produces scalars.
From this single interaction arise norms, orthogonality, distances, and
classical geometric objects such as balls, spheres, and ellipsoids.

\medskip

Many constructions encountered earlier --- bilinear forms, inner products,
matrix multiplication --- share a common feature: they are linear in each
argument separately.
This observation leads naturally to multilinear maps and to tensors.

\medskip

Tensors are treated as abstract multilinear objects defined without
coordinates.
The tensor product is introduced through its universal property, explaining
how multilinear problems are canonically reduced to linear ones.
Index notation, arrays, and computational implementations are presented
explicitly as representations.

\medskip

The text is organized around four successive structural themes: linear
structure, geometric and topological structure, multilinear structure, and
logarithmic structure.
Each theme is a refinement of the same organizing idea: structure is revealed
through invariance under transformations.
The logarithmic theme, in particular, illustrates that this principle is not
confined to the linear setting --- the logarithm is a global isomorphism
between multiplicative and additive structures, and as such it belongs to the
same invariant framework as vectors and tensors.

\medskip

For the reader, in order to understand the subjects and use these notes as a
starting point for further research, prior familiarity is assumed with basic
set theory~\cite{wilder}\cite{cohen}\cite{apostol}, relations, functions
\cite{apostol}, $n$-tuples~\cite{iribarren}\cite{cohen}, and basic
algebra~\cite{iribarren}.
Basic knowledge of calculus~\cite{apostol} is also assumed.
The reports~\cite{kolecki} and~\cite{kolecki2} cover the foundational material
on tensors.
Other introductory material is scattered throughout the bibliography.

\medskip

The topics gathered in this monograph belong to the established literature on
linear algebra and abstract algebra.
What motivated its writing was therefore the perception that a single organizing idea ---
the preservation of structure under transformation --- recurs persistently
across both pure and applied science.
To gather that unity into a single volume is the purpose of this work.
Throughout the text, the word \textit{invariant} is used in a broad sense:
that which remains unchanged when coordinates, representations, or physical
substrates are varied.
The reader acquainted with algebra will recognize that the term also carries
more precise technical meanings --- invariant subspaces, similarity invariants,
topological invariants --- which should not be confused with this structural
and transversal usage.

\medskip

An AI language model was used for editing, \LaTeX{} and stylistic refinement.

\medskip

This monograph presents the abstract structure that underlies every algorithm
and autonomous computing system, and therefore remains invariant under any
change in the state of technology.
There is a more physical and hardware-oriented counterpart, centered on a
1960s computer --- the Apollo Guidance Computer --- which is more specific to a
moment in the history of the art.
For the full context, see:

\medskip

{\small
\url{https://jportillo34.github.io/ApolloGuidanceComputer/}
}

\clearpage

\makeatletter
\@mainmattertrue
\makeatother
\pagenumbering{arabic}

\chapter{Linear Transformations, Matrices, and Composition}

\[
\Big\langle
\big\langle K, +_g, \cdot \big\rangle,\;
\big\langle E, +_E \big\rangle,\;
\circ
\Big\rangle
\]

The triple above defines a \emph{vector space $E$} over the \emph{field $K$} \cite{cohen}. Elements of $K$ are called \emph{scalars}, and elements of $E$ are called \emph{vectors}. The operation
\(
\circ : K \times E \to E
\)
is scalar multiplication, and $+_E$ is vector addition. In what follows, symbols $\alpha_1,\dots,\alpha_k$ always denote scalars in $K$, and symbols $x_1,\dots,x_k$ always denote vectors in $E$.

\medskip

Informally, a vector space is a set of vectors that can be added together and scaled by elements of a field.

\bigskip
\hrule
\medskip

\textit{Example (Coordinate $n$-space, \cite{greub}).}

Let $\Gamma$ be a field. Consider the set
\[
\Gamma^n = \Gamma \times \dots \times \Gamma
\]
of $n$-tuples $(\xi^1,\dots,\xi^n)$, with $\xi^i \in \Gamma$. Define addition by

\[
(\xi^1,\dots,\xi^n)+(\eta^1,\dots,\eta^n)=(\xi^1+\eta^1,\dots,\xi^n+\eta^n)
\]

and scalar multiplication by

\[
\lambda(\xi^1,\dots,\xi^n)=(\lambda\xi^1,\dots,\lambda\xi^n),\quad \lambda \in \Gamma.
\]

With these operations, $\Gamma^n$ is a vector space over $\Gamma$, called the \emph{$n$-space over $\Gamma$}. In particular, $\Gamma$ itself is a vector space over $\Gamma$, where scalar multiplication coincides with field multiplication.

\medskip
\hrule
\bigskip

Given vectors $x_1,\dots,x_k \in E$ and scalars $\alpha_1,\dots,\alpha_k \in K$, an expression of the form
\[
\alpha_1 \circ x_1 +_E \alpha_2 \circ x_2 +_E \dots +_E \alpha_k \circ x_k
\]
is called a \emph{linear combination}.

\medskip

A subset $S \subset E$ is called a \emph{system of generators} (or a \emph{spanning set}) of $E$ if every vector $x \in E$ can be written as a linear combination of elements of $S$ \cite{greub}. That is, for every $x \in E$ there exist vectors $s_1,\dots,s_k \in S$ and scalars $\alpha_1,\dots,\alpha_k \in K$ such that
\[
x = \alpha_1 \circ s_1 +_E \dots +_E \alpha_k \circ s_k.
\]

\medskip

A finite family of vectors $\{x_1,\dots,x_k\} \subset E$ is called \emph{linearly independent} if the relation
\[
\alpha_1 \circ x_1 +_E \dots +_E \alpha_k \circ x_k = 0
\]
implies
\[
\alpha_1 = \dots = \alpha_k = 0.
\]
If such a nontrivial relation exists (i.e., with at least one coefficient $\alpha_i\neq 0$), the family is called \emph{linearly dependent}.

\medskip

A finite set of vectors $\{b_1,\dots,b_n\} \subset E$ is called a \emph{basis} of the vector space $E$ if it is both a system of generators of $E$ and linearly independent.

\medskip

Equivalently, $\{b_1,\dots,b_n\}$ is a basis of $E$ if every vector $x \in E$ can be written \emph{uniquely} as a linear combination
\[
x = \alpha_1 \circ b_1 +_E \alpha_2 \circ b_2 +_E \dots +_E \alpha_n \circ b_n,
\quad \alpha_1,\dots,\alpha_n \in K.
\]

A basis therefore provides a minimal and non-redundant set of building blocks for the entire space. Once a basis is fixed, every vector is completely determined by its coefficients relative to that basis, so we can say the basis generates the space.

\medskip

If the vector space $E$ admits a finite basis, the number of vectors in any basis of $E$ is called the \emph{dimension} of $E$ and is denoted by $\dim E$.

\bigskip
\hrule
\medskip

% CORRECTION 1: The "Remark on coordinates" originally appeared here and used the term
% "linear isomorphism" before it was formally defined. The remark has been moved to
% immediately after the definition of isomorphism, where the term is available to the reader.

Let $E$ and $G$ be vector spaces over the same field $K$. A function
\(
\psi : E \to G
\)
is called a \emph{linear transformation} (also known as \textit{linear map}) if it preserves linear combinations \cite{iribarren}:
\[
\psi(\alpha_1 \circ_E x_1 + \dots + \alpha_k \circ_E x_k)
=
\alpha_1 \circ_G \psi(x_1) + \dots + \alpha_k \circ_G \psi(x_k).
\]

Linear transformations preserve the algebraic structure of vector spaces independently of any choice of basis.

\medskip

\textit{Definition (Isomorphism).}
Let $E$ and $F$ be vector spaces over the same field $K$.
A linear transformation $\varphi : E \to F$ is called an \emph{isomorphism}
if for every $y \in F$ there exists a unique $x \in E$ such that $\varphi(x) = y$,
and $\varphi(x_1) = \varphi(x_2)$ implies $x_1 = x_2$.
In this case, we say that $E$ and $F$ are \emph{isomorphic}.

\medskip

An isomorphism preserves all linear structure:
addition, scalar multiplication, linear independence, dimension,
and the validity of linear relations.
Hence, isomorphic vector spaces are structurally identical.

\bigskip
\hrule
\medskip

% CORRECTION 1 (continued): Remark on coordinates moved here, after isomorphism is defined.

\textit{Remark (On coordinates and isomorphism \cite{halmos}).}

Let $E$ be a finite-dimensional vector space over $K$ with $\dim E = n$.
Once a basis $B=\{b_1,\dots,b_n\}$ is chosen, every vector $x\in E$
is uniquely represented by its coordinate column
\[
[x]_B \in K^n.
\]
This correspondence defines a linear isomorphism
\[
E \longrightarrow K^n.
\]

\medskip

Thus, from the purely algebraic point of view, every $n$-dimensional
vector space over $K$ is structurally indistinguishable from $K^n$.

\medskip

However, this identification depends on the chosen basis.
Different bases produce different coordinate representations.
The intrinsic properties of the space are precisely those
that remain invariant under such changes of reference system.

\medskip

For this reason, we treat vector spaces as abstract entities,
independent of any particular coordinate realization.
Coordinates are tools for computation; structure is the object of study.

\medskip
\hrule
\bigskip

We define the \emph{kernel} of $\psi$ by
\[
    \ker(\psi)=\{x\in E:\psi(x)=0\},
\]
that is, the set of vectors in $E$ that the transformation sends to the zero vector of $G$. We define the \emph{image} of $\psi$ by
\[
\operatorname{Im}(\psi)=\{\psi(x):x\in E\}\subseteq G
\]

that is, the image is the set of values in $G$ that $\psi$ actually attains.

\medskip

If $E$ is finite-dimensional, then
\[
\dim E = \dim \ker(\psi) + \dim \operatorname{Im}(\psi),
\]
a relation known as the rank--nullity theorem.

\medskip

We say that a linear transformation $\psi$ is \emph{injective} if $\psi(x)=\psi(y)\Rightarrow x=y$ (equivalently, $\ker(\psi)=\{0\}$), and \emph{surjective} if for every $g\in G$ there exists $x\in E$ such that $\psi(x)=g$ (equivalently, $\operatorname{Im}(\psi)=G$).

\medskip

Let $B_E=\{e_1,\dots,e_n\}$ and $B_G=\{g_1,\dots,g_m\}$ be bases of $E$ and $G$. Each image $\psi(e_i)$ can be written uniquely as
\(
\psi(e_i) = \alpha_{1i} \circ g_1 +_G \dots +_G \alpha_{mi} \circ g_m.
\)
The scalars $\alpha_{ji}$ form an $m\times n$ rectangular array
\[
A =
\begin{bmatrix}
\alpha_{11} & \alpha_{12} & \dots & \alpha_{1n} \\
\alpha_{21} & \alpha_{22} & \dots & \alpha_{2n} \\
\vdots & \vdots & & \vdots \\
\alpha_{m1} & \alpha_{m2} & \dots & \alpha_{mn}
\end{bmatrix}.
\]
and its abbreviated form
\[
A=(\alpha_{ij}),
\]
called the \emph{matrix of $\psi$ relative to the chosen bases}.

\medskip

Let $A=(\alpha_{ij})$ be an $m \times n$ matrix over a field $K$.

For each $i=1,\dots,m$, the $i$th row
\[
(\alpha_{i1},\dots,\alpha_{in})
\]
is called a \textit{row vector} of $A$.
It is naturally regarded as an element of $K^n$.

For each $j=1,\dots,n$, the $j$th column
\[
(\alpha_{1j},\dots,\alpha_{mj})^T
\]
is called a \textit{column vector} of $A$.
It is naturally regarded as an element of $K^m$.

\medskip

When $A$ represents a linear transformation $\psi:E\to G$
with respect to chosen bases,
the $j$th column of $A$ is precisely the coordinate vector
of $\psi(e_j)$ in the basis of $G$.

\medskip

This fact follows directly from linearity.
Let $B_E=\{e_1,\dots,e_n\}$ be the chosen basis of $E$.
Every vector $x\in E$ admits a unique representation
\[
x=\alpha_1 e_1+\cdots+\alpha_n e_n.
\]
Since $\psi$ preserves linear combinations,
\[
\psi(x)=\alpha_1\psi(e_1)+\cdots+\alpha_n\psi(e_n).
\]

Thus the action of $\psi$ on an arbitrary vector is completely
determined by its values on the basis vectors.
If the coordinates of each $\psi(e_j)$ relative to the basis of $G$
are arranged as the $j$th column of a matrix $A$, then the coordinate
vector of $\psi(x)$ is obtained by forming the same linear combination
of those columns with coefficients $\alpha_1,\dots,\alpha_n$.
In coordinates,
\[
[\psi(x)]_{B_G}=A[x]_{B_E}.
\]
Thus a linear transformation is completely determined by the images
of the basis vectors, and the matrix records precisely these images
in coordinate form.

\medskip

Matrices are concrete representations of linear transformations once bases have been chosen.

\medskip

Linear transformations may also be applied successively.
Understanding how their coordinate representations combine
leads naturally to matrix multiplication.

\medskip

Matrix multiplication is the coordinate expression of \emph{composition of linear maps}.

\medskip

\textit{Derivation (Matrix multiplication from composition).}

\medskip

Let
\[
S : K^p \longrightarrow K^n,
\qquad
T : K^n \longrightarrow K^m
\]
be linear transformations with matrices (relative to chosen bases)
\[
B \in M_{n\times p}(K),
\qquad
A \in M_{m\times n}(K).
\]

For a vector $x \in K^p$, matrix--vector multiplication represents
application of the corresponding linear map:
\[
S(x)=Bx,
\qquad
T(y)=Ay.
\]

The composition $T\circ S : K^p \to K^m$ therefore satisfies
\[
(T\circ S)(x)=T(S(x))=A(Bx).
\]

Since a matrix represents a linear transformation uniquely in coordinates,
there must exist a matrix $C$ such that
\[
(T\circ S)(x)=Cx
\quad \text{for all } x\in K^p.
\]
By definition, this matrix is the product $C=AB$, and it is determined by
the requirement
\[
(AB)x=A(Bx).
\]
Writing this relation in coordinates shows that the entries of $C=AB$
are necessarily given by
\[
(AB)_{ij}
=
\sum_{k=1}^{n} A_{ik} B_{kj},
\qquad
1\le i\le m,\; 1\le j\le p.
\]

Thus the familiar rule for matrix multiplication is the coordinate expression of composition of linear maps.

\medskip

% CORRECTION 2: The change-of-basis section originally reopened the justification
% for matrix-vector multiplication after it had already been established via composition.
% It is now presented as a distinct application — interpreting vectors as states of a
% system — rather than as a second derivation of the same product rule.

In many practical cases, a vector does not merely represent an abstract element of a
vector space, but the \emph{state} of a system expressed relative to a chosen basis.
We will refer to such a choice of basis as a \emph{reference system}.
This framing adds no new algebraic content, but it motivates a natural question:
what happens to the coordinate representation of a state when the reference
system changes?

\medskip

% CORRECTION 3: The introduction of "state" and "reference system" now includes
% an explicit bridge sentence linking the abstract language used previously
% to the applied terminology being introduced here.

Let $E$ be a finite-dimensional vector space and let
\[
B=\{e_1,\dots,e_n\}
\]
be a basis of $E$. Any vector $x\in E$ can be written uniquely as
\[
x = \alpha_1 \circ e_1 +_E \dots +_E \alpha_n \circ e_n.
\]

The scalars $\alpha_1,\dots,\alpha_n$ are the \emph{coordinates} of the state $x$ relative to the reference system $B$. Writing these coordinates as a column vector,

\[
[x]_B =
\begin{bmatrix}
\alpha_1\\
\vdots\\
\alpha_n
\end{bmatrix},
\]
is a notational convenience: it encodes the action of linear combinations relative to the chosen basis.

\medskip

Suppose now that a second reference system
\[
B'=\{e'_1,\dots,e'_n\}
\]
is chosen for the same space $E$. The vector $x$ is the same abstract object, but its coordinates relative to $B'$ form a different column vector $[x]_{B'}$.

\medskip

The passage from $[x]_B$ to $[x]_{B'}$ is not arbitrary. It is determined by a linear transformation
\[
T : K^n \longrightarrow K^n,
\]
whose matrix depends only on the relation between the two bases. Thus,
\[
[x]_{B'} = A\,[x]_B,
\]
where $A$ is the matrix of the change of reference system.
This is a direct instance of the matrix--vector product $[\psi(x)]_{B_G} = A[x]_{B_E}$
derived above: here the ``transformation'' is precisely the relabelling of coordinates
induced by passing from $B$ to $B'$.

\medskip

In coordinates, each row of the matrix represents the coordinate expression of a linear functional on $K^n$ (we will formalize this identification when discussing dual spaces), and matrix--vector multiplication consists precisely in applying each of these functionals to the coordinate vector. Concretely, if $A$ is an $m\times n$ matrix and $x$ is a column vector in $K^n$, then the $i$th component of $Ax$ is the value of the functional given by the $i$th row of $A$ applied to $x$; this is defined only when $x$ has $n$ components, i.e., only when the number of columns of $A$ equals the number of rows of $x$.

\medskip

The requirement that the number of columns of $A$ equal the number of components of $[x]_B$ is not a convention; it expresses the fact that the codomain of one linear map must match the domain of the next.

\medskip

In this sense, column vectors represent states, matrices represent transformations between reference systems, and matrix multiplication represents successive changes of coordinates. The familiar computational rules are therefore consequences of the abstract structure introduced earlier.

\medskip

% CORRECTION 4: The duplicated sentence about nonlinear systems has been removed.
% Only the more informative version is retained.

Linear transformations therefore form the structural backbone of many quantitative models.
Even in systems where nonlinear effects dominate, computations are typically organized
around repeated applications and compositions of linear maps. The linear framework
therefore provides a universal computational language, even when the phenomena being
modeled are not themselves linear.

\medskip

% CORRECTION 5: A bridge paragraph now connects the composition/matrix discussion
% to the dual space, motivating the new concept from what the reader has already seen.

The derivation of matrix multiplication revealed that a row vector acts on a column
vector to produce a scalar: each row of a matrix defines a rule that assigns a number
to any element of $K^n$, and this rule is linear. This suggests that row vectors
and column vectors are fundamentally different kinds of objects. Making this distinction
precise requires introducing the \emph{dual space}.

\medskip

The set of all linear transformations from $E$ into the base field $K$ is denoted by
\[
E^* = L(E,K)
\]
and is called the \emph{dual space} of $E$. An element $x^* \in E^*$ is called a \emph{linear functional}. Thus,
\[
x^* : E \to K
\]
satisfies
\[
x^*(\alpha_1 \circ_E x_1 +_E \dots +_E \alpha_k \circ_E x_k)
=
\alpha_1 x^*(x_1) + \dots + \alpha_k x^*(x_k).
\]

Linear functionals assign scalar values to vectors while preserving linear structure.

\medskip

If $E$ is finite-dimensional with basis $B_E=\{e_1,\dots,e_n\}$, then every $x^*\in E^*$ is uniquely determined by
\[
x^*(e_1),\dots,x^*(e_n).
\]
These scalars form a row vector
\[
[\,x^*(e_1)\;\dots\;x^*(e_n)\,].
\]

In contrast, vectors in $E$ are represented as column vectors. These two objects live in different spaces: $E$ and $E^*$.

\medskip

Given a matrix
\[
A = (\alpha_{ij}) \in M_{m \times n}(K),
\]
its \emph{transpose} is the matrix
\[
A^T = (\alpha_{ji}) \in M_{n \times m}(K),
\]
obtained by interchanging rows and columns.

\medskip

Thus the $i$th row of $A$ becomes the $i$th column of $A^T$,
and the $j$th column of $A$ becomes the $j$th row of $A^T$.

\medskip

The transpose operation exchanges the roles of rows and columns.
In particular, if a matrix represents a linear map
\( \psi : K^n \to K^m \),
then its transpose has the reversed size and may be viewed,
purely at the level of arrays, as interchanging the domain and codomain indices.

Column vectors represent elements of a vector space $E$, whereas row vectors represent linear functionals, elements of the dual space $E^*$. These are algebraically different objects, even though they may have the same number of components.

\medskip

In many situations, however, a vector is used to define a linear functional. This requires additional structure (Chapter II).

A column vector represents an element of $K^n$, a row vector represents an element of $(K^n)^*$, and an $m\times n$ matrix represents a linear map
\[
K^n \xrightarrow{\;\psi\;} K^m.
\]

A product such as
\[
[\,x^*\,]\,A
\]
is defined because it corresponds to the composition
\[
K^n \xrightarrow{\;\psi\;} K^m \xrightarrow{\;x^*\;} K
\]

The dimensions match precisely because the codomain of the first map equals the domain of the second.

\medskip

By contrast, a product of two row vectors would attempt to compose
\[
K^n \to K
\quad\text{with}\quad
K^n \to K,
\]
which is meaningless: the output of the first map does not lie in the domain of the second. The product is therefore undefined by necessity.

\medskip

Matrix dimensions encode compatibility of domains and codomains. One may multiply matrices if and only if the corresponding linear maps can be composed.

\medskip

\textit{Concrete numerical example.}

Let
\[
x^*(x,y)=3x+2y
\]
be a linear functional on $K^2$. Relative to the canonical basis, it is represented by the row vector
\[
[\,3\;\;2\,].
\]

Let
\[
A=
\begin{bmatrix}
2 & 3\\
5 & 4
\end{bmatrix}
\]
be the matrix of a linear transformation
\(
\psi:K^2\to K^2.
\)

The product
\[
[\,3\;\;2\,]A
=
[\,3\;\;2\,]
\begin{bmatrix}
2 & 3\\
5 & 4
\end{bmatrix}
=
[\,19\;\;17\,]
\]
is defined because it represents the composition
\[
K^2 \xrightarrow{\;\psi\;} K^2 \xrightarrow{\;x^*\;} K.
\]

The result is again a row vector, corresponding to the linear functional
\[
(x,y)\longmapsto 19x+17y.
\]

\medskip

By contrast, the product
\[
[\,3\;\;2\,]\,[\,a\;\;b\,]
\]
is undefined because it would attempt to compose two maps of the form
\[
K^2 \to K,
\]
which is impossible: the output of the first map is a scalar, not a vector in $K^2$.

\bigskip

\textit{Conceptual summary.} The structures introduced in this chapter---vector spaces, linear transformations, bases, coordinate systems, and dual spaces---are all facets of the same underlying linear framework. Matrices encode linear transformations relative to a chosen basis, column vectors encode states relative to a reference system, and systems of linear equations are simply coordinate expressions of these transformations asking whether a vector lies in the image of a map. Algorithms such as Gaussian elimination, while essential for computation, do not create new structures: they explore the coordinate-level consequences of the abstract linear relationships already present. Together, these concepts illustrate how abstract algebraic definitions give rise to the familiar computational and geometric tools used in science and engineering.

\section*{Notes and Remarks}

\textit{Note (Why linearity?, \cite{oxford}).}

\textit{An elementary reason for the foundational nature of linear algebra
within mathematics is that, under suitable conditions, functions can be locally approx-
imated by linear functions. This fact, together with the tendency of natural systems to
reside near the minimum of their potential energy, explains the prevalence of linearity
in scientific phenomena.}

\medskip

\textit{For now, we will say informally that a situation exhibits linearity if we can identify a
numerical output and a numerical input such that the former is proportional to the
latter.}

\bigskip
\hrule
\medskip

\textit{Note: Direct Sums and Decomposition of Vector Spaces \cite{halmos}}

\medskip

Given two vector spaces $U$ and $V$ over the same field $K$,
their \emph{direct sum} is the vector space
\[
U \oplus V
=
\{(x,y) \mid x \in U,\; y \in V\},
\]
with addition and scalar multiplication defined componentwise.

\medskip

The subspaces
\[
\{(x,0) \mid x \in U\}
\quad\text{and}\quad
\{(0,y) \mid y \in V\}
\]
may be identified with $U$ and $V$, respectively.
Every element of $U \oplus V$ admits a unique decomposition
\[
(x,y) = (x,0) + (0,y),
\]
so that $U \oplus V$ is obtained by combining two independent components.

\medskip

More generally, a vector space $W$ is said to be the direct sum of subspaces
$U$ and $V$ if every $w \in W$ can be written uniquely as
\[
w = u + v,
\qquad
u \in U,\; v \in V.
\]
In this case one writes $W = U \oplus V$.

\medskip

The notion of direct sum expresses a structural principle:
a space may be decomposed into independent parts whose contributions
combine additively without interaction.

\medskip

This construction should be distinguished from algebraic extensions.
For example, the complex numbers $\mathbb{C}$ form a two-dimensional
vector space over $\mathbb{R}$ and may be identified, as vector spaces,
with $\mathbb{R} \oplus \mathbb{R}$ via
\[
a + ib \;\longleftrightarrow\; (a,b).
\]
However, this identification does not preserve multiplication.
The product of complex numbers couples the two components,
introducing an interaction that is not present in the direct sum.

\medskip

Thus, while some extensions can be viewed as direct sums at the level of
vector spaces, additional algebraic structure may introduce relations
between components that go beyond linear combination.

\medskip

The direct sum provides the fundamental way of assembling independent
linear components.
Later constructions, such as the tensor product, will introduce a different
kind of combination in which components interact multiplicatively rather
than additively.
\medskip
\hrule
\bigskip

% CORRECTION 6: The Notes are now ordered thematically to provide a clearer
% reading thread. The order is:
%   (a) Why linearity? — motivational framing
%   (b) Direct sums — structural decomposition, extends the chapter's algebra
%   (c) Linear systems — coordinate-level reinterpretation, bridges to computation
%   (d) Matrix as a function — foundational definition, anchors the abstract view
%   (e) Generalized matrix product — algebraic abstraction of multiplication
%   (f) Polynomials — infinite-dimensional example, shows linearity beyond coordinates

\bigskip
\hrule
\medskip

\textit{Note (Systems of linear equations).}

Operationally, a system of linear equations of the form
\[
\begin{cases}
a_{11} x_1 + \dots + a_{1n} x_n = b_1\\
\vdots\\
a_{m1} x_1 + \dots + a_{mn} x_n = b_m.
\end{cases}
\]

is treated as a collection of equations to be solved for the unknowns
\(
x_1,\dots,x_n.
\)

\medskip

From the abstract point of view, this system is nothing more than the coordinate expression of a linear transformation.

\medskip

Indeed, let
\[
\psi : K^n \longrightarrow K^m
\]
be the linear transformation whose matrix relative to the canonical bases is
\[
A =
\begin{bmatrix}
a_{11} & \dots & a_{1n}\\
\vdots & & \vdots\\
a_{m1} & \dots & a_{mn}
\end{bmatrix}.
\]

\medskip

Writing

\[
x =
\begin{bmatrix}
x_1\\
\vdots\\
x_n
\end{bmatrix}
\quad\text{and}\quad
b =
\begin{bmatrix}
b_1\\
\vdots\\
b_m
\end{bmatrix},
\]
the system above is equivalently expressed as the single equation
\[
\psi(x) = b,
\quad\text{or, in coordinates,}\quad
A x = b.
\]

So, a system of linear equations specifies a vector
\(
b \in K^m
\)
and asks whether it lies in the image of a linear transformation
\(
\psi : K^n \to K^m.
\)
Questions of existence, uniqueness, or multiplicity of solutions correspond to geometric properties of this map, such as injectivity and surjectivity, and are independent of any particular algorithm used to compute them.

\medskip

Algorithms such as Gaussian elimination are methods for finding a vector
\(
x \in K^n
\)
\textit{that satisfies}
\(
A x = b
\)
relative to a chosen reference system (basis). From the abstract perspective, they are simply procedures for exploring the image and kernel of a linear transformation in coordinates. That is, Gaussian elimination exploits the structure of the linear map encoded by the matrix \(A\) to determine whether a solution exists, whether it is unique, and how to express all possible solutions, all while working entirely in the chosen coordinate system.

\medskip
\hrule
\bigskip

\bigskip
\hrule
\medskip

\textit{Note (Matrix as a function, \cite{iribarren}).}

\medskip

Let $S\neq\varnothing$ be a set and let $m,n\ge 1$.
An \emph{$m\times n$ matrix with entries in $S$} is a function
\[
M:\{1,\dots,m\}\times\{1,\dots,n\}\longrightarrow S.
\]
We write $a_{ij}=M(i,j)$ and abbreviate $M$ by $M=(a_{ij})$ (or $M=(a_{ij})_{m\times n}$) when no confusion can arise.

\medskip

\textit{Specialization.}
When $S=K$ is a field, the set of all $m\times n$ matrices over $K$ is denoted by $M_{m\times n}(K)$.

\medskip
\hrule
\bigskip

\bigskip
\hrule
\medskip

\textit{Note (Generalized matrix product, \cite{knuth3}).}

There is a useful abstraction of matrix multiplication that makes explicit which algebraic properties are really being used. If $A$ is an $m\times n$ matrix and $B$ is an $n\times s$ matrix, and if $\circ$ and $\bullet$ are binary operations such that $\circ$ is associative, one may define the \emph{generalized matrix product}
\[
A\,\genfrac{}{}{0pt}{}{\circ}{\bullet}\,B
\]
to be the $m\times s$ matrix $C$ whose entries are
\[
C_{ij} = (A_{i1}\bullet B_{1j}) \circ (A_{i2}\bullet B_{2j}) \circ \cdots \circ (A_{in}\bullet B_{nj}),
\qquad 1\le i\le m,\; 1\le j\le s.
\]
The ordinary matrix product is obtained when $\circ$ is $+$ and $\bullet$ is $\cdot$.

\medskip
\hrule
\bigskip

\bigskip
\hrule
\medskip

\textit{Example (Polynomials as a free linear construction, \cite{knuth2}).}

Let $K$ be a field. Consider the set
\[
K[x] = \left\{
\sum_{k=0}^{n} a_k x^k
\;\middle|\;
n \in \mathbb{N},\; a_k \in K
\right\},
\]
whose elements are called \emph{polynomials} in the formal symbol $x$ with coefficients in $K$.

\medskip

The symbol $x$ is not assigned a numerical value; it serves only as a generator of formal expressions.
Addition and scalar multiplication are defined coefficientwise:
\[
\left( \sum_{k=0}^{n} a_k x^k \right)
+_{K[x]}
\left( \sum_{k=0}^{m} b_k x^k \right)
=
\sum_{k=0}^{\max(n,m)} (a_k +_K b_k) x^k,
\]
\[
\alpha \circ
\left( \sum_{k=0}^{n} a_k x^k \right)
=
\sum_{k=0}^{n} (\alpha \cdot a_k) x^k.
\]

\medskip

With these operations, $K[x]$ becomes a vector space over $K$.

\medskip

The family
\[
\{ 1, x, x^2, x^3, \dots \}
\]
is linearly independent and generates $K[x]$. Hence it forms a basis of $K[x]$.

\medskip

Every polynomial can therefore be written uniquely as a linear combination
\[
u(x)
=
a_0 \circ 1
+_{K[x]}
a_1 \circ x
+_{K[x]}
\dots
+_{K[x]}
a_n \circ x^n,
\qquad a_k \in K.
\]

\medskip

The vector space $K[x]$ is infinite dimensional, since its basis contains infinitely many elements.

\medskip

Conceptually, $K[x]$ is the vector space freely generated by the formal powers
\[
\{ x^k \}_{k \ge 0}.
\]
No relations exist among these generators other than those required by the axioms of a vector space.

\medskip

Polynomial multiplication introduces an additional operation on $K[x]$,
\[
x^i \cdot x^j = x^{i+j},
\]
which is compatible with the linear structure but is not required for the definition of the vector space itself.

\medskip

Thus, polynomials provide a canonical example of a linear structure generated by formal symbols. The linear framework appears independently of geometry, coordinates, or numerical interpretation.
Multiplication enriches this structure, but linearity is primary.

\medskip
\hrule
\bigskip
\chapter{Geometry from Linear Structure}

In Chapter I, vector spaces were introduced as purely algebraic objects.
A vector was an abstract element of a set equipped with addition and scalar multiplication,
and linear transformations were functions preserving this structure.
No notion of length, angle, distance, or orthogonality was assumed or required.

\medskip

Algebra alone allows us to decide whether vectors are independent,
whether a set spans a space,
or whether a linear transformation is injective or surjective.
However, it does not allow us to ask geometric questions.
In many situations, it is not enough to know \emph{which} vectors are possible;
we also need to know how \emph{large} they are, how \emph{close} they are to one
another, or how \emph{aligned} they are.
These notions are expressed through length, distance, and angle.
Without additional structure, the language of vector spaces does not distinguish
between ``larger'' and ``smaller,'' nor does it allow one to speak of
perpendicularity or of the ``best approximation'' of one vector by others
(as in projections and least squares).
There is no intrinsic way, within a bare vector space, to compare the size of
two vectors or to determine whether they point in perpendicular directions.

\medskip

Geometry begins only when additional structure is imposed.

\medskip

It is useful now to recall how geometric space is commonly represented in
coordinates.
Given sets $A_1,\dots,A_n$, their \emph{Cartesian product}
(also known as \emph{Product Set}) is
\[
A_1\times\cdots\times A_n = \{(x_1,\dots,x_n)\mid x_i\in A_i\}.
\]
In particular, the coordinate space $\mathbb{R}^n$ is the Cartesian product of
$n$ copies of $\mathbb{R}$.
This construction by itself is not yet geometry: it provides a convenient
\emph{model} in which ``points'' are represented by tuples of real numbers.
A geometry emerges only after we specify an additional rule that allows us to
measure and compare such tuples.

\medskip

A central theme of this document is that geometry is not a replacement for
linear structure, but an enrichment of it.
The vector spaces introduced earlier remain unchanged;
what changes is that they are now equipped with a rule that allows vectors to
interact in a way that produces scalar quantities.

\medskip

Informally, this rule takes two vectors and returns a scalar.
From this single operation, notions such as length, angle, and distance emerge
naturally.

\medskip

To make this interaction precise, we introduce the notion of a bilinear map.

\medskip

Let $E$ and $G$ be vector spaces over the same field $K$.
A mapping
\[
B : E \times G \longrightarrow K
\]
is called \emph{bilinear} if it is linear in each argument separately.
That is, for fixed $y \in G$, the map
\[
x \longmapsto B(x,y)
\]
is a linear functional on $E$, and for fixed $x \in E$, the map
\[
y \longmapsto B(x,y)
\]
is a linear functional on $G$.

\medskip

Explicitly, for all $x_1,x_2 \in E$, $y_1,y_2 \in G$, and all scalars
$\alpha \in K$,
\[
B(x_1 + x_2, y) = B(x_1, y) + B(x_2, y), \qquad
B(\alpha x, y) = \alpha B(x, y),
\]
\[
B(x, y_1 + y_2) = B(x, y_1) + B(x, y_2), \qquad
B(x, \alpha y) = \alpha B(x, y).
\]

\medskip

When $E = G$, such a mapping is called a \emph{bilinear form} on $E$.

\medskip

At this level, no geometric interpretation is assumed.
A bilinear form is merely a rule that assigns a scalar to an ordered pair of
vectors, subject to linearity in each argument.
It need not be symmetric, positive, or related to any notion of length or angle.

\medskip

Nevertheless, bilinear forms play a fundamental structural role.
They provide a systematic way to associate vectors with linear functionals
and to construct scalar quantities from pairs of vectors.

\medskip

Only after additional conditions are imposed on a bilinear form do familiar
geometric notions emerge.
These special cases will be introduced shortly.

% CORRECTION 3: The restriction to R is now introduced here, immediately after
% bilinear forms and before any mention of positivity or inner products.
% Previously it appeared much later, after the inner product had already been
% implicitly used.

\medskip

Up to this point, the scalar field $K$ has been arbitrary.
Pure linear structure does not require the ability to compare scalars.
Geometry, however, depends on positivity: expressions such as
$\langle x,x\rangle > 0$ and inequalities of the form $r>0$
presuppose an order relation.
For this reason, we now restrict attention to the real number system
\[
\Big\langle \mathbb{R}, +, \cdot, \le \Big\rangle,
\]
viewed as an ordered field.
It is precisely the order on $\mathbb{R}$ that allows lengths, distances, and
orthogonality to acquire geometric meaning.

\medskip

We now specialize the previously introduced bilinear forms.

\medskip

Let $E$ be a vector space over $\mathbb{R}$.
A \emph{bilinear form} on $E$ is a function
\[
B : E \times E \longrightarrow \mathbb{R}
\]
that is linear in each argument separately.
That is, for all $x,x',y,y' \in E$ and all $\lambda \in \mathbb{R}$,
\[
B(x+x',y)=B(x,y)+B(x',y), \qquad
B(\lambda x,y)=\lambda B(x,y),
\]
and similarly in the second argument.

\medskip

Bilinear forms arise naturally from linear transformations.
If
\[
\psi : E \longrightarrow E^*
\]
is a linear mapping, then the formula
\[
B(x,y) = \psi(x)(y)
\]
defines a bilinear form on $E$.
Conversely, every bilinear form determines such a mapping.
Therefore, bilinear forms encode a controlled way in which vectors may interact
to produce scalars.

\medskip

At this level, no geometry is present.
A bilinear form may be degenerate, asymmetric, or indefinite.
It allows vectors to interact, but it does not yet measure anything.

\medskip

Geometry appears when the bilinear form satisfies additional constraints.

% CORRECTION 1 & 2: The inner product is now defined before any coordinate
% expressions involving it. The section on the transpose, which previously
% interrupted the path from bilinear forms to the inner product, has been
% relocated to after the inner product is established. The coordinate expression
% of the inner product follows the definition, not precedes it.

\medskip

An \emph{inner product} on a vector space $E$ over $\mathbb{R}$
is a bilinear form
\[
\langle \cdot,\cdot \rangle : E \times E \longrightarrow \mathbb{R}
\]
that is symmetric, positive, and non-degenerate.
These properties ensure that
\[
\langle x,x \rangle > 0 \quad \text{for all } x \neq 0.
\]

\medskip

Once such a rule is fixed, vectors acquire measurable properties.
Two vectors may be declared orthogonal.
A vector may be assigned a length.
The distance between two points may be defined in terms of the vectors
connecting them.
None of these concepts exist prior to the introduction of this additional
structure.

\medskip

As in Chapter I, coordinates play no fundamental role.
Geometric quantities must be independent of any particular reference system.
A formula that changes when the basis is changed describes a representation,
not a geometric object.

\medskip

Only after a basis is chosen do familiar coordinate expressions appear.
Sums of squares, dot products, and quadratic equations arise as coordinate-level
manifestations of the underlying structure, not as definitions of it.

\medskip

\textit{Coordinate expression of the inner product.}

Let $(E,\langle\cdot,\cdot\rangle)$ be a finite-dimensional inner product space
and let $\{e_1,\dots,e_n\}$ be an orthonormal basis.
Every vector $x\in E$ can be written uniquely as
\[
x=\sum_{i=1}^n x_i e_i,
\qquad
y=\sum_{i=1}^n y_i e_i.
\]

Using bilinearity and the orthonormality condition
$\langle e_i,e_j\rangle=\delta_{ij}$, we obtain
\[
\langle x,y\rangle
=
\sum_{i=1}^n x_i y_i.
\]

Thus, in an orthonormal basis the abstract inner product is represented by the
familiar \emph{dot product} of coordinate vectors in $\mathbb{R}^n$.

\medskip

This observation explains the role of matrix multiplication.
If a linear transformation $\psi:E\to G$ is represented by a matrix
$A=(a_{ij})$ and a vector $x$ has coordinate vector $x=(x_1,\dots,x_n)^T$,
then the $i$-th component of the product $Ax$ is
\[
(Ax)_i=\sum_{j=1}^n a_{ij}x_j.
\]

Each component is therefore the dot product of the $i$-th row of $A$
with the vector $x$.
Matrix multiplication can thus be understood as the systematic evaluation of
inner products between rows of one matrix and columns of another.

\medskip

In this sense, the familiar algebraic rule for multiplying matrices is not an
arbitrary computational convention: it admits a natural geometric interpretation
as the coordinate expression of interactions measured by an inner product.

\medskip

This point of view explains why objects such as circles, spheres, and ellipsoids
can be described by algebraic equations, yet retain a geometric meaning that is
invariant under changes of coordinates.
The equations depend on the chosen basis; the geometry does not.

\medskip

The goal is therefore to show how geometry emerges from linear structure once a
suitable interaction between vectors is specified.

\medskip

% CORRECTION 2 (continued): The section on the transpose and dual transformation
% is now placed here, after the inner product is defined. This preserves the
% algebraic derivation of the transpose as intrinsic (no inner product needed)
% while allowing the remark about duality becoming internal to appear naturally
% in the context of an established inner product.

Although geometry has begun to emerge through bilinear interaction, the
underlying algebraic mechanism remains the linear transformation.
To understand how geometric structure interacts with linear maps, we examine a
construction naturally associated with every linear transformation.

\medskip

The central object of this work is the linear transformation. Before introducing
geometric notions such as length, angle, and distance, we make explicit a
fundamental construction associated with every linear transformation: its
transpose.

\medskip

Let $E$ and $G$ be vector spaces over the same field $K$, and let
\[
\psi : E \longrightarrow G
\]
be a linear transformation. Recall that the dual spaces $E^*$ and $G^*$ consist
of all linear functionals on $E$ and $G$, respectively.

\medskip

For each $g^* \in G^*$, the composition
\[
g^* \circ \psi : E \longrightarrow K
\]
is a linear functional on $E$. This defines a mapping
\[
\psi^* : G^* \longrightarrow E^*,
\qquad
\psi^*(g^*) = g^* \circ \psi.
\]

\medskip

The mapping $\psi^*$ is linear and is called the \emph{transpose}
(or dual transformation) of $\psi$. The direction of the arrow is reversed:
\[
E \xrightarrow{\;\psi\;} G
\quad\Longrightarrow\quad
G^* \xrightarrow{\;\psi^*\;} E^*.
\]

\medskip

This construction is intrinsic: it depends only on the linear transformation
$\psi$ and does not require any choice of basis, coordinates, or inner product.

\medskip

When bases are chosen for $E$ and $G$, and the corresponding dual bases are
chosen for $E^*$ and $G^*$, the linear transformation $\psi$ is represented by
a matrix $A$, and the dual transformation $\psi^*$ is represented by the
transposed matrix $A^T$.
Thus, the familiar matrix transpose is the coordinate expression of the intrinsic
dual transformation $\psi^*$.

\medskip

The inner product does more than assign lengths.
It establishes a canonical identification between the vector space and its dual.

\medskip

For each vector $x \in E$, define the mapping
\[
\Phi(x) : E \longrightarrow \mathbb{R},
\qquad
\Phi(x)(y) = \langle x,y \rangle.
\]

By bilinearity, $\Phi(x)$ is a linear functional on $E$.
Thus, we obtain a mapping
\[
\Phi : E \longrightarrow E^*.
\]

\medskip

If the inner product is non-degenerate, this mapping is injective.
In finite dimensions, since $\dim E = \dim E^*$, it is therefore an isomorphism.

\medskip

Hence, once an inner product is fixed, every vector determines a unique linear
functional, and every linear functional arises from a unique vector.
The distinction between $E$ and $E^*$ disappears at the structural level.

\medskip

It is precisely this identification that allows the transpose of a linear
transformation to be regarded as an operator on $E$ itself.
When geometry is introduced, duality becomes internal.

\medskip

Once an inner product is fixed, vectors acquire length.
The norm of a vector $x \in E$ is defined by
\[
\|x\| = \sqrt{\langle x,x \rangle}.
\]

\medskip

So, the vector space $E$ introduced in Chapter I is now equipped with additional
geometric structure:
\[
\Big\langle
\langle \mathbb{R},+,\cdot\rangle,\;
\langle E,+\rangle,\;
\circ,\;
\langle\cdot,\cdot\rangle
\Big\rangle.
\]

The inner product is not required for linearity itself, but it allows the
introduction of geometric notions such as length, angle, orthogonality, and
adjoint transformations.

\medskip

\textit{Algebraic expansion of the norm.}

For any vectors $x,y \in E$, the square of the norm of their sum satisfies
\[
\|x+y\|^2
=
\langle x+y, x+y \rangle.
\]
By bilinearity and symmetry of the inner product,
\[
\langle x+y, x+y \rangle
=
\langle x,x \rangle
+
\langle x,y \rangle
+
\langle y,x \rangle
+
\langle y,y \rangle
=
\|x\|^2
+
2\langle x,y \rangle
+
\|y\|^2.
\]

Thus,
\[
\|x+y\|^2
=
\|x\|^2
+
\|y\|^2
+
2\langle x,y \rangle.
\]

This identity holds in every inner product space.
It expresses how the interaction term $\langle x,y \rangle$ measures the
deviation from pure additivity of squared lengths.

\medskip

The inner product also satisfies the inequality
\[
|\langle x,y \rangle|
\le
\|x\|\,\|y\|,
\]
known as the \emph{Cauchy--Schwarz inequality}.
As a consequence,
\[
\|x+y\|^2
\le
\|x\|^2
+
\|y\|^2
+
2\|x\|\,\|y\|
=
(\|x\|+\|y\|)^2.
\]

Taking square roots yields the \emph{triangle inequality}
\[
\|x+y\|
\le
\|x\|
+
\|y\|.
\]

\medskip

The \textit{Pythagorean identity} arises as the special case
$\langle x,y \rangle = 0$,
in which the interaction term vanishes and
\[
\|x+y\|^2
=
\|x\|^2
+
\|y\|^2.
\]

\medskip

Orthogonality also becomes meaningful.
Two vectors $x,y \in E$ are said to be orthogonal if
\[
\langle x,y \rangle = 0.
\]

\medskip

The order of the vectors is irrelevant.
Symmetry of the inner product implies
\[
\langle x,y \rangle = \langle y,x \rangle,
\]
so orthogonality is a mutual relation.

\medskip

The zero vector is orthogonal to every vector in $E$.
Indeed, for any $v \in E$,
\[
\langle 0,v \rangle
=
\langle 0\cdot v, v \rangle
=
0\langle v,v \rangle
=
0.
\]
Moreover, the zero vector is the only vector with this property.

\medskip

Thus, orthogonality is precisely the condition under which the interaction term
vanishes.
In that case,
\[
\|x+y\|^2 = \|x\|^2 + \|y\|^2.
\]
Squared lengths become additive exactly when the vectors do not interact.

\medskip

% CORRECTION 4: The interlude on eigenvectors has been moved to here, after
% the norm, orthogonality, and the adjoint have been established. Previously
% it appeared between the Pythagorean identity and the section on distance,
% interrupting the chain norm -> distance -> topology at its most critical point.

\textit{Interlude: Invariant Directions}

Up to this point, linear transformations have been treated algebraically:
a map $\psi : E \to E$ sends vectors to vectors while preserving linear
structure.
No geometric meaning has yet been attached to this action.

\medskip

Once an inner product is introduced, however, the action of a linear
transformation can be examined geometrically.
Given a vector $x \in E$, the image $\psi(x)$ may differ from $x$ in two
independent ways: its length may change, and its direction may change.

\medskip

For a generic vector, both effects occur simultaneously.
Most vectors are stretched and rotated.

\medskip

A remarkable phenomenon occurs when there exists a nonzero vector $x \in E$
such that
\[
\psi(x) = \lambda x
\]
for some scalar $\lambda$.
In this case, the transformation acts on $x$ by pure scaling.
The direction of $x$ is preserved.

\medskip

Such a vector is called an \emph{eigenvector} of $\psi$, and the scalar
$\lambda$ is its corresponding \emph{eigenvalue}.
Eigenvectors identify directions in the space that are geometrically invariant
under the transformation.

\medskip

This invariance is an intrinsic property of the linear transformation itself.
Changing the basis may alter the representation of $\psi$, but it cannot
destroy or create eigenvectors.

\medskip

Under the identification $E \cong E^*$ induced by the inner product, the
transpose $\psi^*$ corresponds to a linear operator on $E$, called the
\emph{adjoint} of $\psi$.

\medskip

When a linear transformation arises from a symmetric bilinear form (or
equivalently, from a self-adjoint operator), its eigenvectors corresponding
to distinct eigenvalues are orthogonal.
These directions define a preferred coordinate system dictated by the geometry,
not by arbitrary choice.

\medskip

It is in this sense that eigenvectors serve as the principal axes of the
ellipsoids associated with quadratic forms.
Diagonalization is the act of aligning coordinates with the invariant directions
already present in the structure.

\medskip

% CORRECTION 4 (continued): With eigenvectors established here, the section on
% distance and topology now follows without interruption.

\textit{Distance and the emergence of geometry.}

Once an inner product is available on a vector space $E$, it determines a norm,
and from the norm we obtain a notion of distance between vectors,
\[
d(x,y)=\|x-y\|.
\]

At this point, no new structure is introduced: distance is merely a
reformulation of algebraic data already present.

\medskip

This distance allows us to describe proximity. For a vector $x \in E$ and a
real number $r > 0$, we consider the set
\[
B(x,r) = \{\, y \in E \mid d(x,y) < r \,\}.
\]
Such a set is called an \textit{open ball} centered at $x$ with radius $r$.
Geometrically, it consists of all vectors that lie closer to $x$ than the
prescribed threshold.

\medskip

The strict inequality is essential. It guarantees that every point inside the
ball admits further nearby points, so that no boundary is included.
This ``breathing room'' is what ultimately allows us to speak of continuity and
deformation.

\medskip

The collection of all open balls determines which subsets of $E$ should be
called open: a set is open if, around each of its points, it contains some open
ball.
In this way, distance generates a notion of openness without any additional
axioms.

% CORRECTION 5: A transitional sentence now prepares the reader for the
% axiomatic definition of a topological space, connecting it to the geometric
% intuition built just above rather than introducing the axioms abruptly.

\medskip

This notion of openness can be captured by a small set of axioms that
preserve its essential properties while allowing the concept to be extended
far beyond the metric setting.
The resulting structure is what is known as a \textit{topology} on $E$.

\medskip

A \textit{topological space} is an ordered pair $(X, \tau)$, where $X$ is a
set and $\tau$ is a collection of subsets of $X$.
For $\tau$ to be a \textit{topology}, it must satisfy the following axioms:

\begin{enumerate}
    \item $\emptyset \in \tau$ and $X \in \tau$.
    \item The union of any collection of sets in $\tau$ is also in $\tau$.
    \item The intersection of any \textit{finite} number of sets in $\tau$ is
    also in $\tau$.
\end{enumerate}

The elements of $\tau$ are called \textit{open sets}.

\medskip

Thus, topology does not appear here as an independent theory, but as a
structural consequence of the metric induced by the inner product.
Geometry gives rise to distance, distance gives rise to neighborhoods, and
neighborhoods give rise to topological structure.

\medskip

\textit{Geometric interpretation.}

In $\mathbb{R}^n$ equipped with its standard inner product, open balls
correspond to the familiar spheres of Euclidean geometry.
More general norms deform these balls into ellipsoidal shapes, reflecting
anisotropic scaling along different directions of the space.
In all cases, each point
\[
x = (x_1,\dots,x_n)
\]
is a vector of the ambient vector space, and the geometric objects considered
are subsets of that space defined by distance relations among its vectors.

\medskip

Once objects have been represented as vectors in this way, the geometric
structures introduced above---norms, distances, balls, and ellipsoids---become
tools for organizing and comparing them.

\medskip

Their relative positions, measured by a chosen metric, encode similarity.
This metric need not be Euclidean: in practice it may be designed to reflect
domain-specific invariances or even learned from data.
Neighborhoods of meaning are then described by open balls, and learning can be
interpreted as the discovery of a geometric and topological organization
underlying the data.

\medskip

Once a norm has been introduced, certain geometric sets arise as direct
consequences of the structure and require no additional assumptions.

\medskip

Let $(E,\langle\cdot,\cdot\rangle)$ be an inner product space and let
$\|x\|=\sqrt{\langle x,x\rangle}$ be the associated norm.
For a fixed vector $x_0\in E$ and a scalar $r>0$, the \emph{closed ball}
of radius $r$ centered at $x_0$ is defined by
\[
B_r(x_0)=\{\,x\in E \mid \|x-x_0\|\le r\,\}.
\]

The corresponding \emph{sphere} is the boundary of this set,
\[
S_r(x_0)=\{\,x\in E \mid \|x-x_0\|= r\,\}.
\]

\textit{Remark (Coordinate description in $\mathbb{R}^n$):}
When $E=\mathbb{R}^n$ is equipped with its standard inner product and we use
the standard coordinates $x=(x_1,\dots,x_n)$, the same intrinsic definitions
take the familiar ``sum of squares'' form.
For a ball centered at the origin,
\[
B_r(0)=\left\{x\in\mathbb{R}^n\;\middle|\;x_1^2+\cdots+x_n^2\le r^2\right\},
\]
and its boundary sphere is
\[
S_r(0)=\left\{x\in\mathbb{R}^n\;\middle|\;x_1^2+\cdots+x_n^2=r^2\right\}.
\]
This is simply the coordinate expression of
$\|x\|=\sqrt{x_1^2+\cdots+x_n^2}$.

\medskip

These definitions are intrinsic and do not depend on any choice of coordinates
or basis.

\medskip

\textit{Remark (Symmetry):}
In the Euclidean case $E=\mathbb{R}^n$ with its standard inner product, the
ball is invariant under orthogonal transformations.
If $Q\in O(n)$, then $\|Qx\|=\|x\|$, and hence
\[
Q\bigl(B_r(0)\bigr)=B_r(0).
\]
More generally, any isometry of an inner product space carries balls to balls.

\medskip

\textit{Remark (Slicing in $\mathbb{R}^n$):}
In $\mathbb{R}^n$, fixing the last coordinate $x_n=t$ shows that the
cross-section of a radius-$r$ ball by the hyperplane $x_n=t$ is an
$(n-1)$-dimensional ball of radius $\sqrt{r^2-t^2}$.
This observation is the geometric content behind the common
``integrate by slices'' approach to $n$-dimensional volume.

\medskip

\textit{Remark (Volume of Euclidean balls):}
In $\mathbb{R}^n$, the $n$-dimensional volume of a Euclidean ball scales as a
constant times $r^n$:
\[
\operatorname{Vol}_n\bigl(B_r(0)\bigr)=\omega_n\,r^n,
\]
where $\omega_n=\operatorname{Vol}_n\bigl(B_1(0)\bigr)$ is the volume of the
unit ball. A closed form is
\[
\omega_n=\frac{\pi^{n/2}}{\Gamma\!\left(\frac{n}{2}+1\right)}.
\]
In particular, $\omega_1=2$, $\omega_2=\pi$, and $\omega_3=\frac{4\pi}{3}$.

\medskip

\textit{Remark (A high-dimensional effect):}
Although $B_1(0)\subset \mathbb{R}^n$ always contains many points, its
$n$-dimensional volume satisfies $\omega_n\to 0$ as $n\to\infty$.
This illustrates a basic phenomenon of high-dimensional geometry: sets that
look large in low dimensions can occupy a vanishing fraction of ambient volume
as the dimension increases.

\medskip

More generally, let $B:E\times E\to \mathbb{R}$ be a symmetric
positive-definite bilinear form. The set
\[
\mathcal{E}=\{\,x\in E \mid B(x-x_0,x-x_0)\le r^2\,\}
\]
is called an \emph{ellipsoid}.
Thus, an ellipsoid is the unit ball associated with a bilinear form, or
equivalently, with a modified inner product on $E$.

\medskip

When the bilinear form $B$ is symmetric and positive-definite, there exists a
unique self-adjoint linear operator $A:E\to E$ such that
\[
B(x,y)=\langle Ax,y\rangle
\]
for all $x,y\in E$.

\medskip

The operator $A$ encodes the geometry of the ellipsoid.
Since $A$ is self-adjoint on a finite-dimensional inner product space, it
admits an orthonormal basis of eigenvectors.
If $\{e_1,\dots,e_n\}$ is such a basis and $A e_i=\lambda_i e_i$ with
$\lambda_i>0$, then the defining inequality
\[
B(x-x_0,x-x_0)\le r^2
\]
becomes
\[
\sum_{i=1}^n \lambda_i\, \xi_i^2 \le r^2,
\]
where $\xi_i=\langle x-x_0,e_i\rangle$ are the coordinates in this
eigenbasis.

% CORRECTION 6: The observation that diagonalization "reveals invariant
% directions already present" appeared twice in close succession. The first
% occurrence (here, in the ellipsoid derivation) is retained as it concludes
% the algebraic argument. The redundant restatement that appeared a few lines
% later, just before the coordinate formula, has been removed.

\medskip

Thus, the principal axes of the ellipsoid are precisely the eigenvectors of
$A$, and the eigenvalues determine the scaling along those directions.
Diagonalization does not create the axes; it merely reveals the invariant
directions already present in the structure.

\medskip

Accordingly, when a basis is chosen in which the bilinear form is diagonal,
this invariant definition reduces to a familiar coordinate expression.
In $\mathbb{R}^n$, one obtains
\[
\mathcal{E}
=
\left\{
(x_1,\dots,x_n)\in\mathbb{R}^n
\;\middle|\;
\sum_{i=1}^n \frac{(x_i-a_i)^2}{h_i^2} \le r^2
\right\}.
\]

Written componentwise, this inequality reads
\[
\left(\frac{x_1-a_1}{h_1}\right)^2
+
\left(\frac{x_2-a_2}{h_2}\right)^2
+
\cdots
+
\left(\frac{x_n-a_n}{h_n}\right)^2
\le r^2.
\]

The equation depends on the chosen coordinates; the ellipsoid itself does not.

\medskip

\emph{Remark (Degenerate cases).}
The geometric nature of the set defined above depends on the coefficients
$h_1,\dots,h_n$.
If $h_1=\cdots=h_n$, the ellipsoid reduces to an $n$-dimensional sphere.
If one or more of the parameters $h_i$ vanish, the defining quadratic form
becomes degenerate and the ellipsoid collapses into a lower-dimensional
object: a disk, a line segment, or, in the extreme case, a single point.

\medskip

From the structural point of view, these degenerations correspond to a loss of
rank in the associated quadratic or bilinear form.
Thus, changes in the algebraic structure translate directly into changes in
geometric dimension.

\medskip

\emph{Remark.}
Throughout this discussion, the ambient space $E=\mathbb{R}^n$ is a vector
space in the sense of Chapter I.
Each element
\[
x=(x_1,\dots,x_n)
\]
is therefore a vector of $E$.

\medskip

However, the sets defined by these expressions are \emph{subsets} of $E$
selected by quadratic constraints.
They are not vector spaces themselves, since they are not closed under
addition or scalar multiplication.

\bigskip

\textit{Conceptual summary.}
Vector spaces provide the stage.
Linear transformations describe allowable changes of state.
Geometry enters only when vectors are permitted to interact through a
structured bilinear rule that produces scalars.
From this interaction, lengths, angles, distances, and geometric shapes emerge.
The transition from algebra to geometry is therefore not a leap, but a
controlled enrichment of structure.
The book \cite{iribarren2} is a wonderful reference on the topology of metric
spaces.

\section*{Notes and Remarks}

\textit{Outlook.}

In contemporary applications such as representation learning, elements of a
set (words, images, or signals) are mapped into high-dimensional vector spaces,
where geometric relations encode similarity and structural information.
These constructions provide concrete realizations of abstract linear-algebraic
ideas.

\medskip

Modern computational systems offer another illustrative example.
Platforms devoted to \emph{observability}---that is, the systematic measurement
of the internal state of a running system---collect large families of numerical
quantities such as latency, throughput, memory consumption, error rates, and
resource utilization.

\medskip

At a fixed instant of time, these measurements may be assembled into a vector
\[
x(t) = (x_1(t),\dots,x_n(t)) \in \mathbb{R}^n,
\]
whose coordinates represent observable characteristics of the system.
As time evolves, the system determines a trajectory in this ambient linear
space, so that monitoring and analysis may be interpreted geometrically as
the study of paths in a finite-dimensional vector space.

\medskip

Operational considerations typically impose bounds on each observable,
expressed by inequalities of the form
\[
a_i \le x_i \le b_i,
\]
reflecting physical limitations or acceptable operating ranges.
Geometrically, these constraints define a Cartesian product of intervals,
namely an $n$-dimensional hyperrectangle (or, after normalization, a
hypercube).
Normal system behavior may therefore be viewed as confined to a bounded region
of the ambient space, while anomalous behavior corresponds to trajectories
departing from this region.

\medskip

It should be emphasized that this construction is an idealized mathematical
model: in practical systems the collection of observables may vary over time,
so the dimension $n$ is best regarded as fixed only locally or within a chosen
representation.
Nevertheless, the linear-space framework provides a natural language for
describing states, constraints, and evolution.

\medskip

In this way, observability systems illustrate a recurring theme of linear
algebra: a flat ambient space serves as a universal stage upon which meaningful
phenomena are characterized by geometrically defined subsets.
Norm constraints give rise to spheres, anisotropic scalings produce ellipsoids,
and coordinate-wise bounds define hyperrectangles.
Despite their differing geometric appearances, all arise as structured subsets
embedded within the same linear ambient space.

\bigskip
\hrule
\medskip

% CORRECTION 7: The historical remark on Lambert has been moved to immediately
% after the Interlude on Invariant Directions in the Notes section, where its
% connection to eigenvectors and angle preservation is contextually immediate.
% Previously it appeared later in the notes, separated from the discussion that
% motivated it.

\textit{Historical remark (Angles and invariance).}

Before eigenvectors were understood as invariant directions, their detection
was often indirect.
In two dimensions, determining whether a transformation preserves a direction
reduces to detecting whether it preserves an angle.

\medskip

This problem led naturally to the computation of arctangents.
In the eighteenth century, J.H.~Lambert introduced the continued fraction
\[
\arctan(z)
=
\cfrac{z}{1 + \cfrac{z^2}{3 + \cfrac{4z^2}{5 + \cfrac{9z^2}{7 + \cdots}}}},
\]
as part of his work on orbital mechanics and rotational phenomena.

\medskip

From the modern point of view, this analytic effort reflects a geometric
question: does a linear transformation rotate a direction, or does it merely
scale it?
Eigenvectors provide the structural answer to this question, rendering angle
computation auxiliary rather than fundamental.

\medskip
\hrule
\bigskip

\textit{Remark (Quadratic forms and the chi-square statistic).}
Expressions of the form defining ellipsoids appear naturally in contexts that
are not, at first sight, geometric.
A notable example arises in the theory of random number testing, where one
encounters the quantity
\[
V
=
\frac{(Y_1 - np_1)^2}{np_1}
+
\frac{(Y_2 - np_2)^2}{np_2}
+
\cdots
+
\frac{(Y_n - np_n)^2}{np_n},
\]
known in statistics as the \emph{chi-square statistic} associated with the
observed quantities $Y_1,\dots,Y_n$ and reference values $np_1,\dots,np_n$.

\medskip

From the present point of view, this expression is simply a quadratic form on
$\mathbb{R}^n$.
It measures the squared length of the deviation vector
\[
(Y_1 - np_1,\dots,Y_n - np_n)
\]
with respect to a symmetric positive-definite bilinear form whose matrix, in
the standard basis, is diagonal with entries
$(np_1)^{-1},\dots,(np_n)^{-1}$.

\medskip

Consequently, the sets defined by inequalities of the form
\[
V \le c
\]
are ellipsoids centered at the point $(np_1,\dots,np_n)$.
The denominators $np_i$ determine the anisotropic scaling along the coordinate
directions, exactly as in the general ellipsoid
\[
\sum_{i=1}^n \frac{(x_i-a_i)^2}{h_i^2} \le r^2.
\]

\medskip

Thus, before any probabilistic interpretation is imposed, the chi-square
expression is already a geometric object: it is the squared norm induced by a
particular inner product on $\mathbb{R}^n$.
Statistical theory makes use of this geometry, but does not create it.

\medskip

\textit{Remark (The ambient vector space and constrained regions).}
The vector space $E$ itself is the \emph{ambient space} in which all vectors
live.
In coordinates, when a basis is chosen, this space appears as $\mathbb{R}^n$,
a flat linear environment in which vectors may be freely added and scaled.

\medskip

Geometric objects such as balls, spheres, and ellipsoids do not replace this
ambient space; rather, they are subsets of it defined by constraints on the
norm or on a quadratic form.
For example, the sphere
\[
S_r(x_0)=\{x\in E \mid \|x-x_0\|=r\}
\]
selects vectors whose distance from $x_0$ is fixed, while an ellipsoid selects
vectors whose coordinates satisfy an anisotropic quadratic inequality.

\medskip

In many modern applications (AI Deep Learning for instance), algorithms operate
on vectors in a high-dimensional ambient space while imposing such constraints
during computation.
Normalization steps may restrict vectors to spheres, regularization procedures
may confine parameters to ellipsoidal regions, and optimization methods search
for solutions within neighborhoods described by balls.
The ambient vector space provides the linear universe in which vectors exist,
while these geometric subsets describe regions selected by the rules or
constraints of the problem.
\chapter{Tensors as Multilinear Maps}

Linearity expresses that a map respects the algebraic structure of a vector space.
Many constructions in mathematics, however, depend on functions that take several
vector arguments simultaneously while remaining linear in each argument when the
others are held fixed. This leads naturally to the notion of multilinearity.

\medskip

Let \(E_1, E_2, \dots, E_n\) and \(F\) be vector spaces over the same field \(K\).

\begin{definition}[Multilinear map]
A map
\[
T : E_1 \times E_2 \times \cdots \times E_n \to F
\]
is called \emph{multilinear} if it is linear in each argument separately.
That is, for each index \(i \in \{1,\dots,n\}\), and for all vectors
\[
x_i, y_i \in E_i, \quad \alpha \in K,
\]
we have
\[
T(x_1,\dots,x_i +_{E_i} y_i,\dots,x_n)
=
T(x_1,\dots,x_i,\dots,x_n)
+_F
T(x_1,\dots,y_i,\dots,x_n),
\]
and
\[
T(x_1,\dots,\alpha \circ x_i,\dots,x_n)
=
\alpha \circ T(x_1,\dots,x_i,\dots,x_n),
\]
with all other arguments held fixed.
\end{definition}

When \(n=1\), a multilinear map is simply a linear map.
Thus multilinear maps are not a new kind of object, but a direct generalization
of linear transformations.

\medskip

Multilinear maps arise naturally in many contexts:
bilinear products, pairings between vectors and linear functionals,
and higher-order interactions that cannot be reduced to a single input.
Their systematic study leads to the concept of tensors.

\medskip

Informally, a tensor is an object designed to encode multilinear behavior
in an intrinsic and coordinate-free way.

\medskip

There are two equivalent viewpoints that will be used throughout this document.
The first treats tensors directly as multilinear maps.

\medskip

Let \(E\) be a vector space over \(K\), and let \(E^*\) denote its dual space,
that is, the vector space of linear maps from \(E\) to \(K\).

\begin{definition}[Tensor as a multilinear map]
A \emph{tensor of type \((k,\ell)\)} on \(E\) is a multilinear map
\[
T : \underbrace{E^* \times \cdots \times E^*}_{k\ \text{times}}
\times
\underbrace{E \times \cdots \times E}_{\ell\ \text{times}}
\;\longrightarrow\; K.
\]
\end{definition}

Such a tensor takes \(k\) linear functionals and \(\ell\) vectors as inputs,
and produces a scalar.
Linearity is required in each argument separately.

\medskip

This definition makes precise several familiar cases:
\begin{itemize}
\item A scalar in \(K\) can be viewed as a tensor of type \((0,0)\).
\item A vector in \(E\) corresponds to a tensor of type \((0,1)\).
\item A linear functional in \(E^*\) is a tensor of type \((1,0)\).
\item A bilinear form on \(E\) is a tensor of type \((0,2)\).
\item If $E$ is finite-dimensional, a linear map $T : E \to E$ can be identified
      with a tensor of type $(1,1)$.
\end{itemize}

Thus vectors and matrices reappear as the lowest-order cases of a single
unified framework.

\medskip

At this level of abstraction, no coordinates, arrays, or geometric interpretations
are involved.
A tensor is defined entirely by how it acts on vectors and linear functionals,
and by the multilinearity of this action.

% CORRECTION 1: The motivation for the tensor product is now clearly separated
% from its definition. The previous version mixed motivation and solution in the
% same block, with the final sentence anticipating the definition about to follow
% (circular). The motivation is now self-contained: it identifies the limitation
% of multilinear maps (no natural composition) and poses the question the tensor
% product will answer, without naming the solution prematurely.

\medskip

The definition above is conceptually clear, but it has a practical limitation.
Linear maps compose: if $S : E \to F$ and $T : F \to G$ are linear, then
$T \circ S : E \to G$ is again linear, and its matrix is the product of the
matrices of $T$ and $S$.
Multilinear maps do not admit an analogous operation.
Their outputs and inputs do not naturally match in a way that allows canonical
composition, and there is no general rule for combining two multilinear maps
into a third while preserving multilinearity.

\medskip

Multilinear maps do form vector spaces under pointwise operations, so they can
be added and scaled.
What is missing is a way to \emph{reduce} multilinearity to linearity: to
replace a multilinear problem by a linear one on a suitably constructed space.
This is precisely what the tensor product achieves.

\medskip

Let \(E_1,\dots,E_n\) be vector spaces over the same field \(K\).
Consider the set of all multilinear maps
\[
T : E_1 \times \cdots \times E_n \longrightarrow F
\]
into an arbitrary vector space \(F\).
The tensor product is defined so that every such multilinear map corresponds
uniquely to a linear map defined on a single vector space.

\medskip

\begin{definition}[Tensor product]
The \emph{tensor product} of the vector spaces \(E_1,\dots,E_n\) is a vector
space, denoted
\[
E_1 \otimes \cdots \otimes E_n,
\]
together with a multilinear map
\[
\otimes : E_1 \times \cdots \times E_n \longrightarrow
E_1 \otimes \cdots \otimes E_n,
\]
satisfying the following universal property:

For every vector space \(F\) and every multilinear map
\[
T : E_1 \times \cdots \times E_n \longrightarrow F,
\]
there exists a unique linear map
\[
\widetilde{T} : E_1 \otimes \cdots \otimes E_n \longrightarrow F
\]
such that
\[
T(x_1,\dots,x_n) = \widetilde{T}(x_1 \otimes \cdots \otimes x_n)
\]
for all \(x_i \in E_i\).
\end{definition}

\medskip

This property characterizes the tensor product completely.
The elements \(x_1 \otimes \cdots \otimes x_n\) are called \emph{simple tensors}.
Every element of the tensor product space is a finite linear combination of
simple tensors.

\medskip

As a consequence, the study of multilinear maps is equivalent to the study of
linear maps defined on tensor products.
Specifically,
\[
\text{Multilinear maps } E_1 \times \cdots \times E_n \to F
\quad\longleftrightarrow\quad
\text{Linear maps } E_1 \otimes \cdots \otimes E_n \to F.
\]

\medskip

This equivalence is the conceptual foundation of tensor theory.
It explains why tensors can be manipulated algebraically, composed with linear
maps, and represented concretely once bases are chosen.

\medskip

Returning to the case of a single vector space \(E\), a tensor of type
\((k,\ell)\) may equivalently be viewed as a linear map
\[
T :
\underbrace{E^* \otimes \cdots \otimes E^*}_{k}
\otimes
\underbrace{E \otimes \cdots \otimes E}_{\ell}
\;\longrightarrow\; K.
\]

\medskip

Thus tensors are linear maps defined on tensor products, and their apparent
complexity arises only when coordinates are introduced.

\medskip

In contrast with the direct sum (Chapter~I), which combines independent
components additively, the tensor product encodes interactions between
components through multilinearity.

\bigskip
\hrule
\medskip

\textit{Note:} Once bases are chosen, tensor products become finite-dimensional
vector spaces, and tensors admit coordinate representations as multidimensional
arrays. These representations depend on the chosen bases and do not define the
tensor itself.

\medskip
\hrule
\bigskip

Let $E$ be a finite-dimensional vector space over $K$, and let
\[
B = \{ e_1, \dots, e_n \}
\]
be a basis of $E$. Its dual basis
\[
B^* = \{ e^1, \dots, e^n \}
\subset E^*
\]
is defined by the relations
\[
e^i(e_j) = \delta^i_j,
\]
where $\delta^i_j$ is the Kronecker delta.

\medskip

Every vector $x \in E$ admits a unique expansion
\[
x = x^i \circ e_i,
\]
and every linear functional $\varphi \in E^*$ admits a unique expansion
\[
\varphi = \varphi_i \circ e^i.
\]

The scalars $x^i$ and $\varphi_i$ are the \emph{components} of $x$ and
$\varphi$ relative to the chosen bases.

\medskip

Let now $T$ be a tensor of type $(k,\ell)$ on $E$, that is,
\[
T : \underbrace{E^* \times \cdots \times E^*}_{k}
\times
\underbrace{E \times \cdots \times E}_{\ell}
\longrightarrow K.
\]

The action of $T$ is completely determined by its values on basis elements:
\[
T(e^{i_1}, \dots, e^{i_k}, e_{j_1}, \dots, e_{j_\ell})
= T^{i_1 \dots i_k}{}_{j_1 \dots j_\ell}.
\]

These scalars are called the \emph{components of $T$ relative to the basis $B$}.
They encode the tensor entirely once a basis has been fixed.

\medskip

Given arbitrary arguments
\[
(\varphi^1, \dots, \varphi^k, x_1, \dots, x_\ell),
\]
multilinearity implies the expansion
\[
T(\varphi^1, \dots, \varphi^k, x_1, \dots, x_\ell)
=
T^{i_1 \dots i_k}{}_{j_1 \dots j_\ell}
\,
\varphi^1_{i_1} \cdots \varphi^k_{i_k}
\,
x_1^{j_1} \cdots x_\ell^{j_\ell}.
\]

Thus index notation is the coordinate expression of multilinearity after a
basis is chosen.

\medskip

Bilinear maps provide the simplest computational example of multilinearity.
A bilinear map
\[
B : V \times W \to U
\]
is linear in each argument separately. Once bases are chosen in the vector
spaces involved, the map is represented by coefficients $t_{ij}^{\;k}$ such
that the components of the output vector satisfy
\begin{equation}
z^{k} = \sum_{i=1}^{m}\sum_{j=1}^{n} t_{ij}^{\;k}\, x^{i} y^{j}.
\end{equation}

The collection of coefficients $(t_{ij}^{\;k})$ forms a rank-three tensor.
This representation shows that familiar computational procedures are coordinate
expressions of multilinear mappings. For example, the multiplication of complex
numbers
\begin{equation}
(x_1 + i x_2)(y_1 + i y_2)
= (x_1 y_1 - x_2 y_2)
+ i(x_1 y_2 + x_2 y_1)
\end{equation}
can be interpreted as the evaluation of two bilinear forms determined by a
$2 \times 2 \times 2$ tensor of coefficients.

From this viewpoint, algorithms manipulating numerical arrays are
implementations of tensorial transformations written in a chosen reference
system. Computation therefore makes explicit the invariant multilinear
structure underlying algebraic operations.

\medskip

In particular, a bilinear map
\[
B : E \times F \to G
\]
between finite-dimensional vector spaces is represented, once bases are chosen,
by components
\[
B^k{}_{ij},
\]
so that for vectors $x \in E$ and $y \in F$,
\[
(B(x,y))^k = B^k{}_{ij}\, x^i y^j.
\]
Thus multidimensional arrays of coefficients appearing in numerical algorithms
are coordinate representations of tensors encoding multilinear operations.

\medskip

Once a basis has been fixed, the components
\[
T^{i_1 \dots i_k}{}_{j_1 \dots j_\ell}
\]
can be arranged in a multidimensional array of scalars.

\medskip

Changing the basis of $E$ changes the numerical values of the components
according to precise transformation rules, while the tensor itself remains
unchanged.

\bigskip
\hrule
\medskip

\textit{Note: Tensors via Transformation Laws}

The definition of a tensor as a multilinear map is intrinsic and
coordinate-free. An equivalent characterization emphasizes how the components
of a tensor transform under a change of basis.

\medskip

This viewpoint makes explicit the invariant content that survives changes of
coordinates, which is precisely what allows tensorial laws to appear in physics
and in engineering applications independent of the chosen representation.

\medskip

Let $B = \{e_i\}$ and $B' = \{e'_i\}$ be two bases of $E$, related by
\[
e'_i = A^j{}_i \, e_j,
\]
where $A^j{}_i$ is an invertible change-of-basis matrix.
The corresponding dual bases satisfy
\[
e'^i = (A^{-1})^i{}_j \, e^j.
\]

\medskip

Let $T$ be a tensor of type $(k,\ell)$ on $E$, with components relative to $B$
given by
\[
T^{i_1 \dots i_k}{}_{j_1 \dots j_\ell}.
\]

The matrix $A^j{}_i$ expresses the new basis vectors in terms of the old ones,
as discussed in Chapter~I.
Accordingly, the transformation rule below describes how the coordinate
representation of $T$ changes when passing from the basis $B$ to $B'$.

\medskip

The components relative to the new basis $B'$ are given by
\[
T'^{i_1 \dots i_k}{}_{j_1 \dots j_\ell}
=
A^{i_1}{}_{p_1}
\cdots
A^{i_k}{}_{p_k}
\,
(A^{-1})^{q_1}{}_{j_1}
\cdots
(A^{-1})^{q_\ell}{}_{j_\ell}
\,
T^{p_1 \dots p_k}{}_{q_1 \dots q_\ell}.
\]

\medskip

Contravariant indices (upper indices) transform with the matrix $A$, while
covariant indices (lower indices) transform with its inverse.
This transformation rule is not an additional assumption.
It is a direct consequence of multilinearity and of how vectors and dual vectors
transform under a change of basis.

\medskip

Conversely, suppose that for each choice of basis one assigns a multidimensional
array of scalars
\[
T^{i_1 \dots i_k}{}_{j_1 \dots j_\ell}
\]
in such a way that these arrays are related by the transformation law above
whenever the basis is changed.
Then these coordinate collections define a unique multilinear map
\[
T :
\underbrace{E^* \times \cdots \times E^*}_{k}
\times
\underbrace{E \times \cdots \times E}_{\ell}
\longrightarrow K,
\]
independent of the chosen basis.

\medskip

Thus the following two viewpoints are equivalent:
\begin{itemize}
\item A tensor is a multilinear map.
\item A tensor is a collection of components that transform according to the
      tensor transformation law.
\end{itemize}

The first definition emphasizes intrinsic structure.
The second emphasizes invariance under change of reference frame.
Their equivalence explains why tensor theory plays a central role in geometry
and physics: tensorial equations remain valid in every coordinate system,
because both sides transform according to the same law.

\medskip
\hrule
\bigskip

Thus:
\begin{itemize}
\item tensors are abstract multilinear objects;
\item arrays are basis-dependent representations;
\item indices record how multilinearity decomposes into coordinates.
\end{itemize}

This distinction is essential for interpreting computational implementations.

\medskip

Up to this point, tensors have been considered as static multilinear objects.
We now introduce the operation that allows them to interact algebraically.

% CORRECTION 3: The broken sentence ("operation. we now introduce...") and the
% repetition of the no-composition remark have both been removed. The transition
% is now a single clean sentence that opens the contraction section.

\medskip

To combine tensors and produce new multilinear objects, one introduces the
fundamental operation of tensor contraction, which generalizes matrix
multiplication and explains the appearance of summation over repeated indices.

\medskip

We denote by
\[
\mathcal{T}^{(k,\ell)}(E)
\]
the vector space of all tensors of type $(k,\ell)$ on $E$.

\medskip

Let $E$ be a finite-dimensional vector space over $K$.
Consider a tensor
\[
T \in \mathcal{T}^{(k,\ell)}(E)
\]
and a tensor
\[
S \in \mathcal{T}^{(m,n)}(E).
\]

Assume that one covariant slot of $T$ and one contravariant slot of $S$ are
selected.
Tensor contraction is the operation that pairs these two slots using the
natural evaluation map
\[
E^* \times E \to K.
\]

\begin{definition}[Tensor contraction]
Let
\[
T \in \mathcal{T}^{(k,\ell)}(E),
\qquad
S \in \mathcal{T}^{(m,n)}(E).
\]
Choosing one covariant slot of $T$ and one contravariant slot of $S$, their
contraction is defined by composing these arguments with the natural evaluation
pairing
\[
\langle\cdot,\cdot\rangle : E^* \times E \to K.
\]
In coordinates, this operation appears as summation over a basis.

\medskip

The resulting tensor has type
\[
(k+m-1,\;\ell+n-1).
\]
\end{definition}

\medskip

Contraction reduces the total number of tensor slots by two: one contravariant
and one covariant slot are paired and eliminated.
No coordinates are required to define this operation.

\medskip

To express contraction in coordinates, a basis must be chosen.
The abstract pairing then appears as summation over repeated indices.

\medskip

Fix a basis $\{e_i\}$ of $E$ and its dual basis $\{e^i\}$.

Let
\[
T^{i}{}_{j}
\quad\text{and}\quad
S^{j}{}_{k}
\]
be the components of two tensors of type $(1,1)$.

\medskip

The contraction over the index $j$ produces a new tensor with components
\[
R^{i}{}_{k}
=
T^{i}{}_{j} \, S^{j}{}_{k}.
\]

Here the repeated index $j$ indicates summation:
\[
R^{i}{}_{k}
=
\sum_{j=1}^{n}
T^{i}{}_{j} \, S^{j}{}_{k}.
\]
This summation is not an additional rule.
It is the coordinate expression of the abstract pairing between $E^*$ and $E$.

\medskip

Let $A : E \to E$ and $B : E \to E$ be linear maps.
Each can be identified with a tensor of type $(1,1)$.
Relative to a basis, their components are written as
\[
A^{i}{}_{j},
\qquad
B^{j}{}_{k}.
\]

The composition $A \circ B : E \to E$ has components
\[
(A \circ B)^{i}{}_{k}
=
A^{i}{}_{j} \, B^{j}{}_{k}.
\]

Thus matrix multiplication is exactly tensor contraction between two $(1,1)$
tensors, and the familiar summation rule for matrix products is a special case
of the abstract pairing between $E^*$ and $E$.

\medskip

At the abstract level we have:
\[
\text{Tensor contraction}
\]
which, after choosing bases, becomes
\[
\text{Index summation}
\quad\longrightarrow\quad
C^{i}{}_{k} = \sum_j A^{i}{}_{j} B^{j}{}_{k}.
\]

% CORRECTION 4: The paragraph beginning "It is fundamental to recognize..."
% has been removed. Its tone was markedly more rhetorical than the surrounding
% text, it introduced the non-standard term "index demolition", and its content
% (contraction as the engine of computation) is already conveyed more precisely
% by the Systolic Arrays note below.

% CORRECTION 2: The discrete grid example (force array) has been moved to here,
% after contraction is defined. The reader now has the algebraic tools to
% understand what the indices of the array mean and how they could be contracted.
% Previously it appeared before contraction, leaving the array without any
% operational context.

\medskip

In computational and modeling contexts, tensors frequently appear through
multidimensional arrays whose indices label independent parameters of a problem.
The following example illustrates how such arrays arise in practice while still
representing tensorial data.

\medskip

Consider a physical system modeled on a discrete two-dimensional grid.
At each spatial point $(i,j)$ and at each time step $t$, a force vector
\[
F(t,i,j) \in \mathbb{R}^2
\]
is measured, with components $(F_x, F_y)$ relative to a fixed spatial frame.

\medskip

Let:
\begin{itemize}
\item $t \in \{1,\dots,T\}$ index time,
\item $(i,j) \in \{1,\dots,N\} \times \{1,\dots,M\}$ index spatial position,
\item $\alpha \in \{1,2\}$ index force components.
\end{itemize}

The collected data can be written as an array
\[
F_{t i j \alpha},
\]
with shape $(T,N,M,2)$.

\medskip

Each index corresponds to an independent modeling choice: time, space, and
vector components.
The array stores numerical values relative to chosen coordinate systems.
Only the component index $\alpha$ corresponds to a genuine vector slot in the
tensorial sense; the indices $(t,i,j)$ label independent parameters of the
model.
Contracting the index $\alpha$ against a linear functional produces, at each
spacetime point, a scalar measurement of the force along a chosen direction.

\medskip

Abstractly, this data can be interpreted as the coordinate representation of a
tensor-valued function on a discrete domain.

\bigskip
\hrule
\medskip

\textit{Note: Linear Algebra Realized in Hardware --- Systolic Arrays
\cite{kung}}

\medskip

The interpretation of matrix multiplication as tensor contraction is not only
a conceptual reformulation.
It also determines how modern computational hardware is designed.

\medskip

In large numerical computations, the dominant operation is the repeated
evaluation of inner products,
\[
c_i = \sum_{j=1}^{n} a_{ij} b_j,
\]
which is the coordinate expression of a linear map applied to a vector.
Equivalently, it is a contraction between a tensor of type $(1,1)$ and a
tensor of type $(0,1)$.

\medskip

In the late 1970s, H.~T.~Kung and C.~E.~Leiserson proposed organizing
computation around this operation through architectures called
\emph{systolic arrays}.
Instead of moving data back and forth between memory and a central processor,
data flows rhythmically through a grid of simple processing elements.
Each element repeatedly performs a single operation of the form
\[
C \leftarrow A B + C,
\]
a multiply-accumulate step implementing one contribution to an inner product.

\medskip

Mathematically, each processing element evaluates a partial tensor contraction.
As data propagates across the array, successive contractions accumulate, and
the global result emerges from the coordinated action of many local linear
operations.

\medskip

This organization reflects a structural principle: hardware efficiency is
achieved when computation mirrors algebraic structure.
Since linear transformations decompose into sums of products, an architecture
optimized for multiply-accumulate operations naturally implements linear algebra
at scale.

\medskip

Modern high-performance computing architectures frequently adopt variants of
this paradigm, arranging many simple processing elements to evaluate
multiply-accumulate operations in parallel.
What appears at the software level as tensor contraction or matrix
multiplication is therefore executed physically as streams of coordinated inner
products.

\medskip

From the present viewpoint, this development is not accidental.
Tensor contraction is the fundamental interaction between multilinear objects,
and large-scale computation reduces complex algorithms to repeated evaluations
of this interaction.

\medskip

In this sense, numerical arrays stored in memory, index summations written in
formulas, and data flowing through silicon circuits are three realizations of
the same mathematical object: the evaluation of multilinear maps.

\medskip
\hrule
\bigskip

The abstract principles developed above appear concretely in classical physics,
where tensorial quantities describe invariant relations between physical
observables.

\medskip

\textit{A classical example: the inertia tensor \cite{feynman}}

In classical mechanics, the distribution of mass in a rigid body is encoded by
a second-order tensor called the \emph{inertia tensor}.
Relative to a chosen orthonormal basis of $\mathbb{R}^3$, its components are
given by
\[
I_{ij} = \int_{\Omega}
\big( \delta_{ij} \|x\|^2 - x_i x_j \big)\, dm,
\]
where $\Omega$ is the body and $x$ is the position vector.

\medskip

The inertia tensor defines a linear relation between angular velocity $\omega$
and angular momentum $L$:
\[
L_i = I_{ij} \, \omega^j.
\]

\medskip

Under a change of basis determined by an invertible matrix $A$, its components
transform according to
\[
I'_{ij}
=
A^{p}{}_{i} \,
A^{q}{}_{j} \,
I_{pq},
\]
which is precisely the transformation law of a tensor of type $(0,2)$.

\medskip

Thus, although the numerical entries of $I_{ij}$ depend on the chosen reference
frame, the geometric object it represents does not.
The physical content---resistance to rotational acceleration---is invariant
under change of coordinates.

\medskip

Tensor theory therefore extends the central principle introduced in Chapter~I:
mathematical structure is identified not by its numerical representation,
but by the invariance of its defining relations under admissible
transformations.

% CORRECTION 5 & 6: The Notes are now reordered to group applied examples
% together and separate them from the speculative outlook.
% Order: Outlook (future hardware) -> chi-square remark (bridge back to Ch. II)
% -> OLAP cubes (data analysis application) -> Systolic arrays (hardware).
% The redundant introductory paragraph that preceded the Systolic Arrays note
% in the main text has been removed; the note's own opening now suffices.

\section*{Notes and Remarks}

\bigskip
\hrule
\medskip

\textit{Outlook: Computation Beyond Electronics \cite{Xu}}

\medskip

Recent developments in computing architectures explore physical substrates
beyond traditional electronic circuits, including optical and photonic systems
in which information is processed through interference, diffraction, and
propagation of light.

\medskip

Despite their physical differences, such systems are designed to implement the
same fundamental algebraic operations encountered throughout this chapter.
In particular, the dominant computational task remains the evaluation of
expressions of the form
\[
C^{i}{}_{k} = A^{i}{}_{j} B^{j}{}_{k},
\]
that is, tensor contraction.

\medskip

In these emerging architectures, the interaction of physical signals naturally
realizes multiply-accumulate operations in parallel, so that the evaluation of
multilinear maps is performed directly by the underlying physical process rather
than by sequential electronic instructions.

\medskip

A non-trivial constraint, however, arises from the physical encoding of
information.
Non-coherent optical systems represent values as light intensities, which are
intrinsically non-negative; they therefore operate within the restricted
domain~$\mathbb{R}^+$ rather than the full real line.
Since tensor contractions require weights and activations that range over all
of~$\mathbb{R}$, such systems are algebraically incomplete for general
computation.
Coherent architectures address this limitation by encoding values as optical
field \emph{amplitudes}---quantities that carry sign---and recovering the
signed output via interference at detection.
The full group $(\mathbb{R},+)$ is thereby made available, and tensor
contractions of arbitrary sign are realizable in hardware.

\medskip

From this perspective, advances in hardware do not alter the mathematical
structure of computation.
They provide new physical realizations of the same invariant operations: linear
transformations and tensor contractions, now extended to their natural domain.
The algebraic completeness required here---the passage from~$\mathbb{R}^+$
to~$\mathbb{R}$---is precisely the structural distinction between a
multiplicative and an additive group, whose correspondence is the subject of
Chapter~IV.

\medskip

Thus, independently of whether computation is carried out by electronic,
optical, or other physical means, the algebraic framework developed here
remains unchanged.

\medskip
\hrule
\bigskip

\bigskip
\hrule
\medskip

\textit{Remark: Bridge to Chapter~II --- the chi-square statistic as a
$(0,2)$ tensor}

\medskip

The chi-square statistic encountered in Chapter~II provides an earlier example
of a tensor of type $(0,2)$.
Recall that the expression
\[
V
=
\sum_{i=1}^{n}
\frac{(Y_i - np_i)^2}{np_i}
\]
was identified as the squared norm induced by a symmetric positive-definite
bilinear form $B$ on $\mathbb{R}^n$, whose matrix in the standard basis is
diagonal with entries $(np_i)^{-1}$.
From the present point of view, $B$ is precisely a tensor of type $(0,2)$: a
multilinear map
\[
B : \mathbb{R}^n \times \mathbb{R}^n \longrightarrow \mathbb{R}
\]
that is symmetric and positive-definite.
Its coordinate representation changes with the basis, but the geometric
object it defines---the ellipsoidal level sets of $V$---does not.
Thus the statistical quantity $V$, the quadratic form of Chapter~II, and the
$(0,2)$ tensor of the present chapter are three descriptions of the same
invariant structure.

\medskip
\hrule
\bigskip

\bigskip
\hrule
\medskip

\textit{Note: OLAP Cubes and Tensor Contraction}

\medskip

In large-scale data analysis, observations are frequently organized along
multiple independent categorical dimensions such as time, location, system
component, or measurement type.
The resulting structure is commonly called an \emph{OLAP cube} (Online
Analytical Processing cube).

\medskip

Mathematically, such an object assigns numerical values to elements of a
Cartesian product
\[
D_1 \times D_2 \times \cdots \times D_n,
\]
where each set $D_i$ represents one observational dimension.
After enumerating these dimensions, the data can be written as a
multidimensional array
\[
F_{i_1 i_2 \dots i_n},
\]
which is precisely the coordinate representation of a tensor indexed by $n$
independent parameters.

\medskip

In practice, only a small subset of all possible index combinations is observed,
so the array is typically sparse.
This again reflects a distinction emphasized throughout this chapter: arrays
depend on chosen coordinates and storage conventions, while the underlying
multilinear structure is independent of representation.

\medskip

Analytical queries often aggregate data along selected dimensions.
For instance, combining observations while ignoring certain indices produces
quantities of the form
\[
G_{i_1 \dots i_k}
=
\sum_{j_1,\dots,j_m}
F_{i_1 \dots i_k j_1 \dots j_m}.
\]
From the tensorial viewpoint, this operation is a contraction: indices
corresponding to aggregated dimensions are summed over, reducing the order
of the tensor.

\medskip

Thus operations known operationally as filtering, grouping, or aggregation
correspond mathematically to projections of a high-dimensional tensor onto
lower-dimensional subtensors.
The OLAP cube therefore provides a concrete realization of tensorial structure
arising in modern data analysis.

\medskip

This example reinforces a recurring theme: tensor contraction appears wherever
multidimensional relations are summarized through systematic aggregation.

\medskip
\hrule
\bigskip
\chapter{Logarithms as Structure-Preserving Transformations}

In the preceding chapters, invariance under change of coordinates was identified
as the defining characteristic of geometric and tensorial objects.
A vector, a bilinear form, a tensor --- each was defined not by its numerical
components, but by the structural relations it preserves under transformation.
The present chapter observes that this same principle operates outside the
linear setting.
The logarithm is an isomorphism between algebraic structures, converting
multiplication into addition while preserving every relation that matters.
The invariant is not a coordinate, a length, or a tensor component ---
it is the algebraic correspondence itself.

% ------------------------------------------------------------------
%\section*{The Central Idea: A Structure-Preserving Map}
% ------------------------------------------------------------------

\medskip

At the most basic level, a logarithm answers the question:
\emph{what exponent produces a given number?}
Yet historically and mathematically,
logarithms owe their power to something deeper:
they convert multiplication into addition --- and they do so exactly.

\medskip

In Chapter~I, an isomorphism was defined as a structure-preserving mapping
between algebraic systems.
The logarithm provides a precise, global realization of that idea.

\medskip

The mapping
\[
\log:\mathbb{R}^+\to\mathbb{R}
\]
is not merely a bijection; it is an isomorphism between two group structures:

\begin{enumerate}
    \item A multiplicative group $(\mathbb{R}^+,\cdot)$.
    \item An additive group $(\mathbb{R},+)$.
    \item An isomorphism
    \[
    \log:(\mathbb{R}^+,\cdot)\longrightarrow(\mathbb{R},+),
    \]
    whose inverse is $\exp:(\mathbb{R},+)\to(\mathbb{R}^+,\cdot)$.
\end{enumerate}

\medskip

The defining property is:
\[
\log(ab)=\log(a)+\log(b),
\qquad
\log(1)=0.
\]

Multiplication corresponds to addition; multiplicative identity corresponds to
additive identity; multiplicative inverses correspond to additive inverses.
No algebraic relation is lost --- only its appearance changes.

\medskip

This is not a numerical trick or a computational shortcut.
It is a structural equivalence: the two groups are, from an algebraic standpoint,
the same object wearing different notation.

% ------------------------------------------------------------------
%\section*{Physical Manifestation: The Slide Rule}
% ------------------------------------------------------------------

\medskip

This isomorphism is made physical in the slide rule.
Distances on a logarithmic scale represent logarithmic values.
Adding distances corresponds to multiplying numbers.

\medskip

The construction of logarithmic scales proceeds by mapping
\[
x \mapsto \log_{10}(x),
\]
placing numbers at distances proportional to their logarithms.
Equal ratios correspond to equal distances.
The structural identity $\log(ab)=\log(a)+\log(b)$ becomes a mechanical operation:
slide one rule against another, and multiplication is performed by addition of lengths.

\medskip

Lengths are not numbers, but they obey the same structural rules.
By exploiting this correspondence, the slide rule becomes an analog computer ---
a physical model of an algebraic isomorphism.

% ------------------------------------------------------------------
%\section*{Logarithms as Area: The Geometric Interpretation}
% ------------------------------------------------------------------

\medskip

Logarithms also admit a purely geometric definition independent of exponentiation.
By construction,
\[
\ln(a)=\int_1^a \frac{1}{x}\,dx,
\]
the signed area under the curve $y=\frac{1}{x}$ from $1$ to $a$.

\medskip

This integral definition makes the isomorphism property transparent.
For $a,b>0$,
\[
\ln(ab)
=\int_1^{ab}\frac{1}{x}\,dx
=\int_1^{a}\frac{1}{x}\,dx+\int_a^{ab}\frac{1}{x}\,dx.
\]
The substitution $x=at$ in the second integral gives
\[
\int_a^{ab}\frac{1}{x}\,dx=\int_1^{b}\frac{1}{t}\,dt=\ln(b),
\]
and therefore $\ln(ab)=\ln(a)+\ln(b)$.
The group isomorphism is a theorem, not an assumption.

\medskip

The integral definition also yields the fundamental derivative formula immediately:
\[
\frac{d}{da}\ln(a)=\frac{1}{a}.
\]

% ------------------------------------------------------------------
%\section*{Rate of Change and Linear Approximation}
% ------------------------------------------------------------------

\medskip

The behavior of $\frac{d}{da}\ln(a) = \frac{1}{a}$ connects the logarithm to
the general framework of derivatives and local linearization.

\medskip

Recall that for a function $y = f(x)$, the derivative at a point $a$ is
\[
f'(a) = \lim_{h\to0}\frac{f(a+h)-f(a)}{h},
\]
the instantaneous rate of change, and geometrically the slope of the tangent line.
Near $a$, a differentiable function behaves approximately like its tangent line:
\[
f(x) \approx f(a) + f'(a)(x-a).
\]

\medskip

This local linearity mirrors the passage from tensors to their components
discussed in Chapter~II: global nonlinear structure becomes tractable when
approximated by its best linear model at each point.
The derivative replaces a nonlinear map by a linear one locally;
the logarithm replaces a multiplicative structure by an additive one globally.
Both are instances of the same organizing principle.

% ------------------------------------------------------------------
%\section*{Numerical Computation: Newton--Raphson Applied to the Logarithm}
% ------------------------------------------------------------------

\medskip

The derivative of the logarithm also enables its efficient numerical computation
via the Newton--Raphson method.
Suppose we wish to find $\ln(x)$, i.e., to solve the equation
\[
e^y = x,
\]
or equivalently $f(y) = e^y - x = 0$.

\medskip

The Newton--Raphson iteration replaces the nonlinear equation by a sequence of
linear ones.
Starting from a guess $y_n$, we approximate $f$ by its tangent line and solve:
\[
0 = f(y_n) + f'(y_n)(y_{n+1}-y_n).
\]
Since $f'(y) = e^y$, this yields
\[
y_{n+1} = y_n - \frac{e^{y_n}-x}{e^{y_n}}
         = y_n + \frac{x - e^{y_n}}{e^{y_n}}.
\]

Starting from $y_0 = 1.5$, successive iterations converge rapidly to
$\ln(5) \approx 1.60944$.

\medskip

The simplicity of this iteration is not accidental.
Working in base $e$ ensures that $\frac{d}{dy}e^y = e^y$,
minimizing algebraic overhead and improving numerical stability.
Logarithms in other bases are recovered by the change-of-base formula:
\[
\log_{b}(x) = \frac{\ln(x)}{\ln(b)}.
\]

% ------------------------------------------------------------------
%\section*{Historical Algorithm: Briggs and Binary Accumulation}
% ------------------------------------------------------------------

\medskip

The same structural idea underlies historical algorithms for constructing
logarithm tables.
Henry Briggs computed base-10 tables by accumulating small multiplicative
corrections; a closely related method suited to binary arithmetic is described
in~\cite{knuth}.

\medskip

The algorithm builds a target number by successive multiplications of the form
\[
\frac{10^k}{10^k-1},
\]
accumulating the corresponding logarithms whenever the partial product does not
overshoot the target.
What appears as computation is, in fact, structure preservation:
each multiplicative step corresponds to an additive increment on the logarithmic
side, faithfully enacting the isomorphism $\log(ab) = \log(a) + \log(b)$.

% ------------------------------------------------------------------
%\section*{The Constant $e$ and Its Representations}
% ------------------------------------------------------------------

\medskip

The natural base $e$ is not an arbitrary choice.
It is the unique base for which $\frac{d}{dx}e^x = e^x$,
and the unique solution to $\int_1^e \frac{1}{x}\,dx = 1$.
Its arithmetic character is revealed by several independent constructions.

\medskip

One classical definition arises from compound interest:
\[
e = \lim_{n\to\infty}\left(1+\frac{1}{n}\right)^n.
\]

Another is provided by its continued fraction expansion:
\[
e = [2;\,1,2,1,1,4,1,1,6,\dots],
\]
written in expanded form as
\[
e
=
2 + \cfrac{1}{1 + \cfrac{1}{2 + \cfrac{1}{1 + \cfrac{1}{1 + \cfrac{1}{4 + \cdots}}}}}.
\]

These representations converge rapidly and are well suited for computation,
but their deeper significance is structural:
$e$ is the natural unit of the isomorphism between $(\mathbb{R}^+,\cdot)$
and $(\mathbb{R},+)$.

% ------------------------------------------------------------------
%\section*{Two Complementary Linearizations}
% ------------------------------------------------------------------

\medskip

The logarithm transforms nonlinear growth into linear motion.
It converts multiplication into addition, products into sums,
and exponential scales into linear ones.

\medskip

Two complementary manifestations of the invariant spirit are therefore visible.
The logarithm is a \emph{global} isomorphism:
it establishes a perfect structural equivalence between two algebraic worlds,
valid everywhere on $\mathbb{R}^+$.
The derivative, by contrast, provides a \emph{local} linearization:
it approximates an arbitrary transformation by its best linear model at each point.

\medskip

In both cases, complexity is rendered intelligible by expressing it through
linear relations.
This is why logarithms appear simultaneously in analysis, geometry, numerical
methods, and physical instruments.
They do not belong to a single domain.
They preserve structure across domains.

\backmatter
\clearpage
\phantomsection
\addcontentsline{toc}{chapter}{Bibliography}

\begin{thebibliography}{99}

\bibitem{iribarren}
Iribarren Ignacio L., \textit{Álgebra Lineal}, Editorial Equinoccio, 2009.

\bibitem{halmos}
Halmos, Paul R., \textit{Finite-Dimensional Vector Spaces}, Springer-Verlag, 1987.

\bibitem{iribarren2}
Iribarren, Ignacio L., \textit{Topología de Espacios Métricos}, Editorial Limusa-Wiley, 1973.

\bibitem{cohen}
Cohen, Leon and Ehrlich, Gertrude, \textit{The Structure of the Real Number System}, D. Van Nostrand Company, 1963.

\bibitem{wilder}
Wilder, Raymond L., \textit{Introduction to The Foundations of Mathematics}, Dover Publications, Inc., 2012.

\bibitem{greub}
Greub, Werner, \textit{Linear Algebra}, Springer-Verlag, 1975.

\bibitem{oxford}
Lukas, Andre, \textit{The Oxford Linear Algebra for Scientists}, Oxford University Press, 2022.

\bibitem{apostol}
Apostol, Tom M., \textit{Calculus Volume I}, John Wiley \& Sons Inc., 1967.

\bibitem{knuth}
Knuth, Donald E., \textit{The Art of Computer Programming Volume I: Fundamental Algorithms}, Addison-Wesley, 1997.

\bibitem{knuth2}
Knuth, Donald E., \textit{The Art of Computer Programming Volume II: Seminumerical Algorithms}, Addison-Wesley, 1981.

\bibitem{knuth3}
Knuth, Donald E., \textit{The Art of Computer Programming Volume I, Fascicle 1: MMIX—A RISC Computer for the New Millennium}, Addison-Wesley, 2005.

\bibitem{feynman}
Feynman Richard P., Leighton Robert, Sands Matthew, \textit{The Feynman Lectures on Physics, Volume II}, California Institute of Technology, 2010.

\bibitem{kolecki}
Kolecki, Joseph C., \textit{An Introduction to Tensors for Students of Physics and Engineering}, NASA/TM-2002-211716, 2002.

\bibitem{kolecki2}
Kolecki, Joseph C., \textit{Foundations of Tensor Analysis for Students of Physics and Engineering With an Introduction
to the Theory of Relativity}, NASA/TP—2005-213115, 2005.

\bibitem{kung}
Kung, H. T., Liserson, Charles L., \textit{Introduction to VLSI Systems - Algorithms for VLSI Processor Arrays}, ed. C. Mead and L. Conway, Addison-Wesley, Reading, MA, 1980.

\bibitem{Xu}
Shaofu Xu, Jing Wang, Haowen Shu, Zhike Zhang, Sicheng Yi, Bowen Bai, XingjunWang, Jianguo Liu, Weiwen Zou, \textit{Optical coherent dot-product chip for sophisticated deep learning regression}, Nature - Light: Science \& Applications, 2021.

\end{thebibliography}

\end{document}
